\documentclass[11pt,reqno]{amsart}
\usepackage[T1]{fontenc}
\usepackage{lmodern,amsmath,amssymb,amsthm,mathtools,mathrsfs}
\usepackage[a4paper,margin=27mm]{geometry}
\usepackage{microtype,booktabs,longtable,array,needspace,etoolbox,enumitem}
\usepackage{xcolor}
\definecolor{MathLink}{HTML}{1F4E79}
\definecolor{SectionLink}{HTML}{00656D}
\definecolor{EquationLink}{HTML}{643F7A}
\definecolor{CitationLink}{HTML}{2F6B3C}
\definecolor{URLLink}{HTML}{843C39}
\usepackage{hyperref}
\hypersetup{colorlinks=true,linkcolor=MathLink,citecolor=CitationLink,
 urlcolor=URLLink,linktoc=all,pdfborder={0 0 0},
 pdftitle={Bedford--Taylor products with a positive closed background},
 pdfauthor={},pdfsubject={Internal research manuscript, 16 November 2026}}
\usepackage[nameinlink,noabbrev]{cleveref}
\newtheorem{theorem}{Theorem}[section]
\newtheorem{proposition}[theorem]{Proposition}
\newtheorem{lemma}[theorem]{Lemma}
\newtheorem{corollary}[theorem]{Corollary}
\theoremstyle{definition}
\newtheorem{definition}[theorem]{Definition}
\newtheorem{convention}[theorem]{Convention}
\theoremstyle{remark}
\newtheorem{remark}[theorem]{Remark}
\newtheorem{example}[theorem]{Example}
\numberwithin{equation}{section}
\newcommand{\Cb}{\mathcal C_b}
\newcommand{\an}{\mathrm{an}}
\newcommand{\supp}{\operatorname{supp}}
\newcommand{\PSH}{\operatorname{PSH}}
\newcommand{\PH}{\mathscr{PH}}
\newcommand{\DSH}{\mathscr{DSH}_{\mathrm{lb}}}
\newcommand{\Mone}{\mathscr M^1}
\newcommand{\Ccone}{\mathscr C}
\newcommand{\ddc}{d'd''}
\newcommand{\Real}{\mathbf R}
\newcommand{\MathSecref}[1]{\begingroup\hypersetup{linkcolor=SectionLink}\hyperref[#1]{Section~\ref*{#1}}\endgroup}
\newcommand{\MathEqref}[1]{\begingroup\hypersetup{linkcolor=EquationLink}\hyperref[#1]{(\ref*{#1})}\endgroup}
\let\OriginalEqref\eqref
\renewcommand{\eqref}[1]{\begingroup\hypersetup{linkcolor=EquationLink}\OriginalEqref{#1}\endgroup}
\newcommand{\prooflink}[1]{\par\noindent\begingroup\hypersetup{linkcolor=SectionLink}\hyperref[#1]{Proof in Section~\ref*{#1}.}\endgroup\par\medskip}
\newcommand{\KeepStatement}[1]{\BeforeBeginEnvironment{#1}{\Needspace{5\baselineskip}}}
\forcsvlist{\KeepStatement}{theorem,proposition,lemma,corollary,definition,convention,remark,example}
\AddToHook{cmd/section/before}{\Needspace{8\baselineskip}}
\AddToHook{cmd/subsection/before}{\Needspace{6\baselineskip}}
\setlength{\emergencystretch}{2em}
\setlist[enumerate]{leftmargin=*,label=\textup{(\roman*)}}
\title[Bedford--Taylor products and relative energy]{Bedford--Taylor products with a positive closed background:\\curvature quotients, relative energy and source equations}
\author{}
\date{16 November 2026\\Internal research manuscript}
\begin{document}
\begin{abstract}
We study Bedford--Taylor products in the original Chambert--Loir--Ducros
calculus over a fixed finite extension of $\mathbf Q_p$. A symmetric
middle-positive current, separately closed for $d'$ and $d''$, is retained
as the background. Locally bounded classical plurisubharmonic potentials
give intrinsic mixed currents with weighted joint decreasing continuity.
The kernel of curvature on bounded classical differences yields a
potential-independent operation, and continuous coefficient relations
extend it to a module of curvature currents. An order estimate then gives
a faithful completion of that module inside the original currents.
On strict Laurent domains whose closed boxes carry the background, the
relative first-energy calculus admits a complete Borel realization with
critical weighted continuity. We describe its variational source map and
prove that a single geometric potential detects the entire original-test
defect of a literal Radon datum. In remaining degree one, finite signed
energy-diffuse data have a unique relative renormalized zero-boundary
solution. The regularization and analytic energy estimates used in the
constructions are stated with their precise hypotheses and cited to the
accompanying foundational manuscripts.
\end{abstract}
\maketitle
\setcounter{tocdepth}{1}
\begingroup\hypersetup{linkcolor=SectionLink}\tableofcontents\endgroup
\clearpage
% BEGIN sections/01_introduction.tex
\section{Introduction}\label{sec:introduction}

The Bedford--Taylor recursion differentiates a potential after multiplying
it by a positive closed current. Over a non-Archimedean field, this
construction raises two distinct questions. The first is analytic:
which potential classes make the recursive products positive and stable
under approximation? The second concerns the objects on which the
construction depends: when can a potential be replaced by another
potential with the same curvature, or a finite curvature expression by
another expression representing the same current?

We develop these questions with one positive closed current retained as
the background. The resulting local calculus leads to an ordered module
of curvature currents. On specified analytic domains it also leads to
relative capacities, a complete first-energy space, and source equations.
The proofs use the local regularization of classical potentials and the
analytic energy estimates in the accompanying manuscripts
\cite{Continuous,Bounded,Order,Energy,Dual}. The arguments below give the
constructions and the transitions between these theories; the cited
analytic inputs are stated at the points where they enter.

\subsection{The local product}

Fix a finite extension $K/\mathbf Q_p$ with its prescribed absolute value.
Let $X$ be a good, Hausdorff, boundaryless, quasi-smooth $K$-analytic
space of pure dimension $n$. Opens are ordinary opens in the Berkovich
topology. On such an open $U$, write $\Cb(U)$ for the actual real-valued,
locally bounded, upper-semicontinuous functions whose pullback to every
analytic curve after every complete isometric extension of $K$ is
subharmonic. These are the locally bounded classical plurisubharmonic
functions of \cite{GR25}. Put $D=d'd''$.

Forms and currents belong to the original Chambert--Loir--Ducros
calculus, with its canonical calibration. Positivity means the middle
positive cone of \cite[Definition 5.1.1(c)]{CLD12}; a positive current is
nonnegative on the corresponding positive complementary forms. A
background of degree $(p,p)$ is a symmetric positive current $S$ such
that $d'S=d''S=0$. All current identities below are tested on every
original compactly supported smooth complementary form. The tests
retain all their actual finite moment refinements, including those on
which the moment map loses rank.

\begin{theorem}[Bounded products with a fixed background]
\label[theorem]{thm:main-bounded}
Let $S$ be a background of degree $(p,p)$ on $U$, and put $q=n-p$.
For $0\le k\le q$ there are unique recursively defined coefficient
Radon currents
\begin{equation}\label{eq:main-recursion}
 T_{S,0}=S,\qquad
 T_{S,k}(u_1,\ldots,u_k)
       =D\bigl(u_1T_{S,k-1}(u_2,\ldots,u_k)\bigr),
       \qquad u_i\in\Cb(U).
\end{equation}
The multiplication in this formula uses the given Borel functions and
the coefficient measures on the analytic space. The currents are
positive, symmetric as currents, and separately $d'$- and $d''$-closed.
The operation is invariant under permutations of the differentiated slots,
additive and positively
homogeneous in each such slot, local for ordinary restriction, and
independent of the regularization used to compute them.

If $u_{i,j}\downarrow u_i$ in $\Cb(U)$ for $0\le i\le k$, with locally
bounded limits, then
\begin{equation}\label{eq:main-decreasing}
 u_{0,j}T_{S,k}(u_{1,j},\ldots,u_{k,j})
       \longrightarrow u_0T_{S,k}(u_1,\ldots,u_k)
\end{equation}
on every original compact smooth complementary test. In top degree this
is local weak convergence of signed Radon measures against compactly
supported continuous functions. Finally, if $D[u_i]=D[v_i]$ for every
differentiated slot, then
\[
 T_{S,k}(u_1,\ldots,u_k)=T_{S,k}(v_1,\ldots,v_k).
\]
\end{theorem}
\prooflink{sec:weighted-continuity}

Coefficient multiplication and the ordered construction are established
in \MathSecref{sec:coefficients-construction}. The convergence proof
uses a common collar on a small norm box. The collar permits Green
exchange with a smooth argument before any two discontinuous
approximation parameters are passed to a limit. Positivity then gives
a lower bound for a fixed trace test. A separate upper bound for
decreasing weights determines the trace limit and hence all original
test limits. The final curvature-independence assertion is proved in
\MathSecref{sec:curvature-quotient}.

\subsection{Curvature relations and order completion}

Let $\Mone(U)$ be the sheaf of coefficient Radon $(1,1)$-currents
locally admitting a finite expression
\begin{equation}\label{eq:main-module}
 A=\sum_\alpha f_\alpha D[g_\alpha],\qquad
 f_\alpha\in C(U,\Real),\quad g_\alpha\in\Cb(U).
\end{equation}
Actual equality of currents is the equivalence relation. Its positive
cone is inherited from the original middle-positive cone. Given
$Q\in\Mone(U)$ positive, define
\[
 \begin{aligned}
 E_Q&=\{A\in\Mone(U):-cQ\le A\le cQ\text{ for some }c\ge0\},\\
 \|A\|_Q&=\inf\{c\ge0:-cQ\le A\le cQ\}.
 \end{aligned}
\]
We denote by $\Mone_{\rm ord}$ the sheaf obtained from the completions
of these normed spaces, using the common control $Q+R$ to compare
limits originally controlled by $Q$ and $R$.

\begin{theorem}[Continuous coefficient action and its order completion]
\label[theorem]{thm:main-module}
For every background $S$ of degree $(p,p)$ and $r\le n-p$, the formula
\begin{equation}\label{eq:main-module-action}
 \Phi_{S,r}(A_1,\ldots,A_r)
 =\sum_{\alpha_1,\ldots,\alpha_r}
       \left(\prod_{i=1}^r f_{i,\alpha_i}\right)
       T_{S,r}(g_{1,\alpha_1},\ldots,g_{r,\alpha_r})
\end{equation}
is independent of all presentations \MathEqref{eq:main-module}. It
defines a symmetric $C(U,\Real)$-multilinear local operation, positive
when every slot is positive.

Each completion $\widehat E_Q$ embeds faithfully and isometrically into
the actual symmetric coefficient Radon currents. These embeddings are
compatible under change of control. The operation $\Phi_{S,r}$ extends
uniquely to $\Mone_{\rm ord}$, locally and independently of the
approximating sequences. If
$-M_iQ_i\le A_i,B_i\le M_iQ_i$ and
$-\delta_iQ_i\le A_i-B_i\le\delta_iQ_i$, then
\begin{equation}\label{eq:main-order-estimate}
 -E R_S\le
 \Phi_{S,r}(A_1,\ldots,A_r)-\Phi_{S,r}(B_1,\ldots,B_r)
 \le E R_S,
\end{equation}
where
\[
 R_S=\Phi_{S,r}(Q_1,\ldots,Q_r),\qquad
 E=\sum_{i=1}^r\delta_i\prod_{j\ne i}M_j.
\]
The output is separately closed whenever it is obtained as a limit of
separately closed mixed approximants.
\end{theorem}
\prooflink{sec:continuous-module}

The bounded-difference criterion is the link between potential theory
and this order structure: for $a,b\in\Cb$, positivity of $D[a-b]$
makes the actual difference $a-b$ classical psh. Freezing continuous
coefficients on a finite cover then transfers their relations to every
fixed background. The multilinear estimate
\MathEqref{eq:main-order-estimate} follows by a positive expansion and
telescoping. The embedding and extension are proved in
\MathSecref{sec:order-completion}.

The continuous coefficient operation acts on degree-one presentations.
A higher-degree analytic relation involving cancellation between several
nonzero mixed terms is a separate relation problem. For example,
potential independence concerns each equality $D[u]=D[v]$, whereas a
quadratic relation can have the form
$\sum_\nu c_\nu D_\delta h_\nu\wedge D_\delta k_\nu=0$.
The latter is the subject of the companion research program.

\subsection{Relative energy and source equations}

The global energy theory uses a domain with an explicit boundary.
Let $B\subset X$ be a compact affinoid, and let
\begin{equation}\label{eq:main-box}
 K_0=\{|f_a|\le R_a,\ |g_b|\ge r_b\}
       \Subset\operatorname{Int}_X(B),\qquad
 V=\{|f_a|<R_a,\ |g_b|>r_b\}\ne\varnothing.
\end{equation}
All thresholds are positive, the analytic functions are defined on $B$,
and $S$ is defined on an ordinary neighborhood of $K_0$.
Put $q=n-p\ge1$ and $d=q+1$.

For a Borel set $E\subset V$, define
\[
 C_S(E;V)=\sup_{\substack{h\in\Cb(V)\\0\le h\le1}}
                  T_{S,q}(h,\ldots,h)(E),\qquad
 C_S^*(E;V)=\inf_{\substack{G\supset E\\G\ {\rm ordinary\ open}}}
                         C_S(G;V).
\]
The actual class $\mathcal E^1_{S,\partial}(V)$ consists of negative
extended classical potentials with compact negative sublevel closures,
canonical truncation no loss, and finite first energy. The precise
definition and the nonpolar construction occur in
\MathSecref{sec:relative-energy}. For its members write
\[
 B(u_0,\ldots,u_q)=\int_V(-u_0)\,D u_1\wedge\cdots\wedge D u_q\wedge S,
 \qquad e(u)=B(u^{[d]}),
\]
and
\[
 I(u,v)=e(u)+e(v)-B(u,v^{[q]})-B(v,u^{[q]}).
\]
A row is strongly Cauchy if
$\sup_{j,l\ge N}I(u_j,u_l)\to0$. Define $\Ccone$ to consist of its
actual Borel limits in every local truncated outer $S$-capacity,
modulo local outer-capacity zero.

\begin{theorem}[The complete relative first-energy space]
\label[theorem]{thm:main-energy}
Under \MathEqref{eq:main-box}, every strongly Cauchy row in
$\mathcal E^1_{S,\partial}(V)$ has a unique limit in $\Ccone$.
There are intrinsic mixed curvature, first-order and energy currents,
and global mixed energies
on $\Ccone$, independent of all defining rows and outer-capacity-null
representatives. The mixed energy is symmetric and satisfies
\begin{equation}\label{eq:main-holder}
 \widehat B(u_0,\ldots,u_q)
       \le\prod_{i=0}^q\widehat e(u_i)^{1/d}.
\end{equation}
The completed discrepancy $\widehat I$ separates relative classes and
is complete. Moreover,
\begin{equation}\label{eq:main-strong-topology}
 \widehat I(u_j,u)\to0
 \quad\Longleftrightarrow\quad
 \begin{cases}
 u_j\to u\text{ in every local truncated outer capacity},\\
 \widehat e(u_j)\to\widehat e(u).
 \end{cases}
\end{equation}
For a finite tuple, let $F$ be continuous on the nonpositive orthant and
satisfy $|F(x)|\le A(1-\sum_i x_i)^a$, with $a\ge0$. Give a
curvature current with $k$ curvature slots the degree budget $b=k$;
for a first-order or energy current over $k$ curvature slots put
$b=k+1$ or $b=k+2$, respectively. In the legal bidegrees, if
$a+b\le q+1$, these multipliers are coefficient-integrable and their
weighted currents converge on every original compact smooth test under
strong convergence of the tuple. The global mixed weighted top measures
also converge. The actual classical class is strongly dense.
\end{theorem}
\prooflink{sec:strong-completion}

The proof first extracts uniform deep tails from the Cauchy condition.
These tails give the uniform integrability required at the critical
degree. Compact sublevels of one fixed late member of each row also
give global tightness, without a common lower barrier. Truncation then
constructs the currents before unbounded multiplication is used.

Put $r(u)=\widehat e(u)^{1/d}$. The completed cone gives the
energy-difference space $\mathbb D_S=\Ccone-\Ccone$ with norm
\[
 \|h\|_{\rm ed}=\inf_{h=a-b}\bigl(r(a)+r(b)\bigr),
          \qquad a,b\in\Ccone.
\]
Write $\mathbb D_{S,\mathrm s}^*$ for its continuous linear
functionals whose restriction to $\Ccone$ is continuous under strong
convergence. A positive Radon measure $\mu$ defines a \emph{literal}
member of this class if
$\ell_\mu(h)=-\int h\,d\mu$ is absolutely defined on $\mathbb D_S$,
independent of outer-null representatives, norm continuous and strongly
continuous on $\Ccone$.

\begin{theorem}[Variational sources and one geometric probe]
\label[theorem]{thm:main-source}
Under \MathEqref{eq:main-box}, $\mathbb D_S$ is Banach and
\[
 J(u)(h)=\widehat B(h,u^{[q]})=-\int h\,d\widehat\mu_u
\]
defines an injective map with norm-closed image in $\mathbb D_S^*$.
The metric $\|J(u)-J(v)\|_*$ is complete and induces the strong topology
on $\Ccone$. Every $\ell\in\mathbb D_{S,\mathrm s}^*$ has a unique
variational maximizer
\[
 u=\Pi(\ell)=\operatorname*{argmax}_{v\in\Ccone}
                     \bigl(\ell(v)-\widehat e(v)/d\bigr)
\]
and a unique decomposition $\ell=J(u)+n$ with
$n(v)\le0$ for every $v\in\Ccone$ and $n(u)=0$.

There is one bounded continuous classical $-1<\psi_0<0$,
determined by the Laurent box and having finite top mass, such that
$\psi_0\pm\epsilon_\phi\phi\in\mathcal E^1_{S,\partial}(V)$
for every original compact smooth real scalar $\phi$ and some
$\epsilon_\phi>0$. The choice of $\psi_0$ and $\epsilon_\phi$ works
for every extending allowed background. For a literal positive Radon
datum as above, set $u=\Pi(\ell_\mu)$ and
\[
 \Xi_{\psi_0}(\mu)=J(u)(\psi_0)-\ell_\mu(\psi_0)\ge0.
\]
Then
\[
 \Xi_{\psi_0}(\mu)=0
 \quad\Longleftrightarrow\quad
 \widehat\mu_u=\mu
 \quad\Longleftrightarrow\quad
 \mu=\widehat\mu_v\text{ for some }v\in\Ccone.
\]
The solution is unique. For every such datum and every original compact
smooth scalar test,
\[
 \left|\int\phi\,d(\mu-\widehat\mu_u)\right|
                    \le\epsilon_\phi^{-1}\Xi_{\psi_0}(\mu).
\]
\end{theorem}
\prooflink{sec:variational-source}

The probe is a convergent sum of exact-moment compensators. Its selected
summands admit both signs of a prescribed compact smooth direction.
The normal inequality on these two competitors gives the defect bound;
literal integration then turns equality on all tests into equality of
Radon measures.

The remaining degree $q=1$ gives a linear equation even when the
ambient dimension is larger than one. For this assertion choose one
compact affinoid $B$, ordinary opens $V\Subset\Omega$ with
$\overline\Omega\subset\operatorname{Int}_X(B)$, and a background
$S$ of degree $(n-1,n-1)$ on $\Omega$. On $F=V\cap\supp S$, the
original compact smooth core has a faithful, regular, strongly local,
coercive Dirichlet-form realization $(\mathcal E,H)$ in a finite
reference measure $\lambda$. This construction and its identification
with the original current are stated in
\MathSecref{sec:linear-dirichlet}. Let $c_V$ be its energy capacity and
$T_k(s)=\max(-k,\min(k,s))$.

\begin{theorem}[The finite diffuse relative Dirichlet equation]
\label[theorem]{thm:main-linear}
In the preceding domain, let $\mu$ be a finite signed Radon measure on
$F$ whose variation vanishes on every $c_V$-polar Borel set. There is a
unique $u\in L^1(F,\lambda)$ with a finite energy-quasicontinuous
representative, $T_ku\in H$ for every $k>0$, the original chain identity
\[
 D(\Phi(u)S)=\Phi'(u)\mu+\Phi''(u)\Gamma_S(u)
\]
for every $\Phi\in C^2(\mathbf R)$ with compactly supported first
derivative, and
$k^{-1}\mathcal E(T_{2k}u-T_ku)\to0$.
Here $\Gamma_S(u)$ is the compatible sigma-finite energy measure of
its clippings. This solution satisfies $D(uS)=\mu$ on every original
compact scalar test and is independent of the faithful reference
measure. If $\mu_j\to\mu$ in total variation, then
\[
 \|u_{\mu_j}-u_\mu\|_{L^1(\lambda)}\to0,
 \qquad
 \mathcal E(T_ku_{\mu_j}-T_ku_\mu)\to0
                       \quad(k>0).
\]
\end{theorem}
\prooflink{sec:linear-dirichlet}

\subsection{Relation to the existing calculus}

The recursive construction and continuous-potential theory originate in
\cite[\S5.6]{CLD12}. The revised treatment
\cite[Lemma 8.5.5 and Propositions 8.7.2, 8.7.4]{CLD25} supplies
positivity, local uniform stability and symmetry for continuous psh
potentials on a positive closed background. Its discussion
\cite[\S8.6.1--8.6.2]{CLD25} isolates the question of annihilating
zero-curvature potential differences against a background. The
regularization manuscripts \cite{Continuous,Bounded} supply the
particular continuous and locally bounded classical approximation
statements used here. The bounded curvature quotient and ordered
difference criterion are developed in \cite{Quotient,Order}.

The relative energy construction in \cite{Energy} works with the
original coefficient measures and the full local test sheaf. Its
global boundary class is tied to \MathEqref{eq:main-box}. The dual
construction in \cite{Dual} then treats a source as a literal measure
when measure equality is required. This separation of local products,
relative completion and source compatibility determines the organization
of the present article.


Wu \cite[Theorem 4.2 and Appendix A]{Wu26} constructs top-degree local
Monge--Amp\`ere measures for a class defined by decreasing tropical-convex
PL approximation on a fixed neighborhood. The classical regularization
input here supplies that neighborhood and approximation from the
complete-extension curve definition; the background calculus retains
all allowed bidegrees. On balanced polyhedral and tropical spaces,
\cite{BMPS26,APW26} develop related product, energy and source theories.
Their geometric categories provide useful comparison for the relation
and solvability questions arising from original analytic currents.

Relative capacity with a positive closed current is classical in the
complex setting \cite{DE06}. Relative non-pluripolar products and their
closedness are developed in \cite[Theorem 3.7]{Vu21}, and the
constant-one first-energy mixed inequality has the complex precedent
\cite[Theorem 3.2]{Cegrell98}. The role of these constructions here is
their realization with the fixed original non-Archimedean background,
coefficient measures and boundary class.

General renormalized measure-data equations for transient regular
Dirichlet forms satisfying the strong sector condition are treated in
\cite{KR15}. In the linear part of this article, the analytic work is
the faithful realization of the part form, its measures and its scalar
chain identities in the original CLD calculus. The abstract
finite-energy nest theorem used there is \cite[Lemma 2.1]{KRStructure}.

% END sections/01_introduction.tex

% BEGIN sections/02_coefficients_and_construction.tex
\section{Coefficient measures and the ordered construction}
\label{sec:coefficients-construction}

\subsection{Positive currents and analytic-space coefficients}

A current is of coefficient Radon type if its coefficient functionals
in every finite original moment presentation are signed Radon measures
on the analytic domain, compatibly with scalar multiplication, relations
between forms, and common refinements. In particular, for every compact
original test $\eta$ there is a signed measure $\mu_{T,\eta}$ satisfying
\[
 \int a\,d\mu_{T,\eta}=\langle T,a\eta\rangle
 \quad(a\text{ original smooth}).
\]
This condition permits the definition
\begin{equation}\label{eq:borel-multiplication}
 \langle hT,\eta\rangle=\int h\,d\mu_{T,\eta}
\end{equation}
for an actual locally bounded Borel function $h$.

\begin{lemma}[Finite test control]\label[lemma]{lem:finite-control}
Every positive current is of coefficient Radon type. For each compact
original test $\eta$ there is a nonnegative compact trace test $\Theta$
and a constant $c_\eta$ such that
\begin{equation}\label{eq:finite-control}
 |\langle T,h\eta\rangle|
       \le c_\eta\|h\|_\infty\langle T,\Theta\rangle
\end{equation}
for positive $T$ and bounded Borel $h$ on the chosen compact stage.
The positive cone is weakly closed and pointed. If positive currents
$T_j$ converge on all original tests, their coefficient variations are
uniformly bounded on each compact stage. Thus
$f_jT_j\to fT$ for locally uniformly convergent continuous $f_j\to f$.
\end{lemma}

\begin{proof}
Localize to a finite moment presentation containing the coefficients of
$\eta$ and a smooth cutoff $\chi$ equal to one near its support.
A positive current annihilates the antisymmetric part of a test.
For the symmetric part, put
\[
 \beta=\sum_{a=1}^N d't_a\wedge d''t_a,
 \qquad\Omega=\beta^m/m!,
\]
where $(m,m)$ is the test bidegree. The coefficient matrix of $\Omega$
is positive definite on this finite exterior space. Matrix domination
therefore gives
$c_\eta\|h\|_\infty\chi\Omega\pm h\eta\ge0$ for smooth $h$.
The inequality persists after adjoining scalar moments and pulling
back to a common refinement. Positivity proves
\MathEqref{eq:finite-control} for smooth $h$.

Compact uniform density of original smooth scalar functions and the
Riesz representation theorem produce the coefficient measures. Their
uniqueness preserves the original relations and refinements. Bounded
Borel approximation against these measures gives
\MathEqref{eq:finite-control} in the stated form. If $T_j\to T$, the
trace masses $\langle T_j,\Theta\rangle$ are bounded, so the same
coefficient estimate is uniform. Approximate $f$ uniformly by a smooth
function, pass the fixed smooth pairing, and then let the uniform error
go to zero; the term $(f_j-f)T_j$ is controlled by the same bound.
Closedness of the positive cone follows on each positive test. If both
$T$ and $-T$ are positive, the trace control makes every pairing zero.
This is the coefficient argument of
\cite[Lemma \emph{Continuous coefficients}]{Continuous} in the present analytic category, using
the ordinary local transport of \cite{Bounded}.
\end{proof}

\subsection{Regularization on a fixed neighborhood}

We use the following two precise regularization consequences.
The first is \cite[Theorem 1.5]{Bounded}; its construction uses the
prescribed field and the original calibration.

\begin{proposition}[Classical regularization]\label[proposition]{prop:regularization}
For every $u\in\Cb(U)$ and every $x\in U$ there is one ordinary
neighborhood $W$ of $x$ and continuous semipositive rational PL
functions $p_j$ over $K$, defined on $W$, with $p_j\downarrow u$.
The class $\Cb$ is stable under decreasing locally bounded limits.
For finitely many decreasing families with locally bounded endpoints,
the common projective extension construction gives one such neighborhood
for all rows. Continuous members of these rows admit smooth positive
uniform approximants on smaller fixed neighborhoods. Smooth majorants
of each bounded member can consequently be chosen pointwise convergent
and locally uniformly bounded.
\end{proposition}

The simultaneous neighborhood assertion uses the norm-box projective
extension in \cite[Appendix D]{Bounded}. Smooth approximation of the
continuous rows is \cite[Lemma \emph{Smooth approximation of FS metrics} and Proposition \emph{Regularized maximum}]{Continuous}.
The smooth approximating rows in the last sentence are not required to
decrease. Their precise collar version is proved below.

\begin{proposition}[First order and ordered products]
\label[proposition]{prop:ordered-construction}
If $R$ is a background and $u\in\Cb$, then $D(uR)$ is positive,
symmetric, separately closed and of coefficient Radon type. Iterating
this operation defines \MathEqref{eq:main-recursion}. The ordered
products are local and separately additive and positively homogeneous
in the potentials. They are independent of the approximants used in
the first-order construction.
\end{proposition}

\begin{proof}
Work on a regularization neighborhood and choose continuous positive
$p_j\downarrow u$. The continuous background theorem
\cite[Theorem \emph{Continuous potentials on a positive closed background}]{Continuous}, transported to this neighborhood as in
\cite{Bounded}, makes $D(p_jR)$ positive. Bounded convergence against
the fixed coefficient measures of $R$ gives $p_jR\to uR$ on every
original test. Distributional differentiation gives
$D(p_jR)\to D(uR)$. Positivity and symmetry pass to the limit.
The identities $d'D=d''D=0$ give separate closedness, and
\cref{lem:finite-control} gives coefficient Radon type. The result is
the derivative of the already defined actual product $uR$, so it is
independent of the row. The same argument applies successively to every
positive closed intermediate current. Linearity of multiplication and
differentiation gives the algebraic properties; restriction commutes
with both operations.
\end{proof}

\begin{proposition}[Relative local mass estimate]\label[proposition]{prop:relative-cln}
Fix compact test stages $L\Subset W$ and $k\le q$. For each finite
original coefficient presentation on $L$ there is a constant $C$ and a
compact positive trace test $\Theta$ on $W$ such that
\[
 \|T_{S,k}(u_1,\ldots,u_k)\|_{L}
       \le C\prod_{i=1}^k\|u_i\|_{L^\infty(W)}
                         \langle S,\Theta\rangle
\]
whenever the displayed bounds are finite. Here the left side denotes
the sum of the coefficient variations in that finite presentation.
The estimate is local and may be applied after every actual refinement.
\end{proposition}

\begin{proof}
Choose nested compactly supported cutoffs between $L$ and $W$ and add
the finitely many moment coordinates used by their coefficients.
The trace estimate in \cref{lem:finite-control} reduces a coefficient
variation to a positive trace pairing. For one differentiated slot,
distributional integration by parts writes it as
$\langle u_1R,D(\chi\alpha)\rangle$, where $\alpha$ is a fixed
closed trace form. On its compact support both signs of
$D(\chi\alpha)$ are bounded by a constant times the next positive
trace form. The coefficient bound is therefore at most
$C\|u_1\|_\infty$ times the trace mass of $R$ on the next compact
stage. Repeat with the remaining differentiated slots. The last trace
is a fixed compact test against $S$. This proves the estimate; finite
localization supplies the general compact stage.
\end{proof}

% END sections/02_coefficients_and_construction.tex

% BEGIN sections/03_weighted_continuity.tex
\section{Weighted decreasing continuity}
\label{sec:weighted-continuity}

The ordered products from \cref{prop:ordered-construction} are defined
before symmetry or joint convergence is used. The following collar
construction supplies the Green exchanges needed to prove convergence.

\begin{lemma}[A common domain for smooth majorants]
\label[lemma]{lem:common-majorants}
Let $x\in U$, and let $u_{i,j}\downarrow u_i$ be finitely many decreasing
families in $\Cb(U)$ with locally bounded limits. After shrinking around
$x$, there is one ordinary neighborhood $W$, independent of $i,j$, with
the following property. Each of the functions $a=u_{i,j}$ and $a=u_i$
has smooth psh majorants $a_\nu$ on $W$ such that
\[
 a_\nu\geq a,\qquad a_\nu(x')\longrightarrow a(x')\quad(x'\in W).
\]
The majorants are bounded uniformly in $\nu$ on each fixed compact subset
after a finite prefix is discarded. If $a\leq M$ on a compact neighborhood
$C\Subset W$, its majorants may be chosen with $a_\nu\leq M+1$ on $C$.
The index discarded may depend on the function $a$.
\end{lemma}

\begin{proof}
Use \cite[Proposition \emph{Transport to an ordinary analytic chart}]{Bounded} to work on one ordinary algebraic
neighborhood, with the original field and calibration. Apply \cite[Lemma D.2, \emph{A common collar for decreasing families}]{Bounded} to the entire finite family of sequences, before choosing
approximants. It gives one projective variety, one ample line, one frame
on $W$, and bounded classical metrics whose weights on $W$ are the given
functions plus one fixed constant. By \cite[Theorem \emph{Projective contact and regularization}]{Bounded}, every such
metric is a decreasing limit of pure Fubini--Study metrics on this same
line. Their weights $p_\nu$ are continuous, satisfy $p_\nu\geq a$, and
decrease pointwise to $a$ after the fixed constant is removed.

Smooth each pure Fubini--Study metric uniformly, using \cite[Lemma \emph{Smooth approximation of FS metrics} and
Proposition \emph{Regularized maximum}]{Continuous}. If the uniform error is at most
$\epsilon_\nu\downarrow0$, add $\epsilon_\nu$ to the smoothed weight.
The resulting smooth psh function $a_\nu$ satisfies
\[
 p_\nu\leq a_\nu\leq p_\nu+2\epsilon_\nu.
\]
This proves pointwise convergence and the majorant property.
On $C$, the increasing open sets $\{p_\nu<M+1/2\}$ cover $C$ whenever
$a\leq M$. Compactness gives a common index on $C$. Discard this prefix
and require $2\epsilon_\nu<1/2$. The same argument with any bound for
$a$ gives local uniform boundedness of the row. No monotonicity of the
\emph{smoothed} row is asserted or needed.
\end{proof}

\begin{lemma}[A smooth psh collar]
\label[lemma]{lem:smooth-collar}
Under the assumptions of \cref{lem:common-majorants}, there are
compact polynomial norm boxes
\[
 x\in\operatorname{Int}C_1,\qquad
 C_1\subset\operatorname{Int}C_{\mathrm{mid}},\qquad
 C_{\mathrm{mid}}\subset\operatorname{Int}C_2\Subset W,
\]
a smooth psh function $b$ on a neighborhood of $C_2$, and a nonnegative
original smooth cutoff $\chi$ compactly supported in
$\Omega=\operatorname{Int}C_2$, with the following properties.
After adding one constant $c$ to all the original functions, put
\[
 v_{i,j}=\max\{u_{i,j}+c,b\},\qquad
 v_i=\max\{u_i+c,b\}.
\]
Then:
\begin{enumerate}
\item $v_{i,j},v_i\in\Cb(\Omega)$ and $v_{i,j}\downarrow v_i$;
\item $v_{i,j}=u_{i,j}+c$ and $v_i=u_i+c$ on a neighborhood of $C_1$;
\item $\chi=1$ on a neighborhood of $C_{\mathrm{mid}}$;
\item all $v_{i,j},v_i$ equal $b$ on one fixed neighborhood of
      $\operatorname{supp}d'\chi\cup\operatorname{supp}d''\chi$;
\item each patched function $v$ has smooth psh majorants $v_\nu\geq v$
      converging pointwise to $v$, uniformly bounded on $\operatorname{supp}\chi$,
      and equal to $b$ on that same collar.
\end{enumerate}
The majorant rows in (5) need not be decreasing.
\end{lemma}

\begin{proof}
Choose the boxes by the construction in \cite[Lemma D.2, \emph{A common collar for decreasing families}]{Bounded}. In
particular, lower-bound polynomials have no zero near $C_2$. Choose
thresholds strictly between the defining thresholds for $C_1$ and
$C_{\mathrm{mid}}$. The finite maximum
\[
 b_0=\max\{0,\log|P_a^{N_a}/c_a|,
                     \log|d_l/Q_l^{M_l}|\}
\]
is zero on a neighborhood of $C_1$ and at least some $\delta>0$ on a
neighborhood in $C_2$ of $C_2\setminus\operatorname{Int}C_{\mathrm{mid}}$.
All its branches other than the zero branch are logarithms of analytic
units where they are active. A branch with a zero is strictly inactive
near that zero.

Apply the regularized maximum of \cite[Proposition \emph{Regularized maximum}]{Continuous}, with
smoothing radius smaller than all the fixed gaps. The result $\widetilde b$
is an original smooth psh function: zero branches at polynomial zeros
are deleted by the exact inactivity property. On a neighborhood of
$C_1$, only the constant zero branch is active, so $\widetilde b=0$.
On the outer collar, $\widetilde b\geq\delta$.

The endpoint bounds give a constant $c$ and a finite $M$ such that
\[
 1\leq u_i+c\leq u_{i,j}+c\leq M
 \quad\hbox{on a neighborhood of }C_2.
\]
Multiply $\widetilde b$ by a positive constant so that the resulting $b$
is greater than $M+3$ on the outer collar. Choose $\chi$, with
$0\leq\chi\leq1$, equal to one near $C_{\mathrm{mid}}$ and compactly
supported in $\Omega$. Smooth cutoff existence and a finite localization
are part of the original CLD test calculus. Its derivative support lies
where the fixed outer-collar inequality holds.

Classical psh functions are stable under maxima and addition of constants.
This proves (1)--(4). To prove (5), apply \cref{lem:common-majorants}
to the \emph{unpatched} function $a=u_{i,j}+c$ or $u_i+c$. Discard a
prefix so that its smooth psh majorants $a_\nu$ are at most $M+1$ on
$C_2$. Define
\[
 v_\nu=\operatorname{rmax}_{\epsilon_\nu}(a_\nu,b),
 \qquad \epsilon_\nu\downarrow0,
\]
with all radii less than the fixed collar gap. The regularized maximum
majorizes the ordinary maximum and has error at most $\epsilon_\nu$.
It is exactly $b$ on the collar because the other branch is inactive
there. These facts give all the assertions in (5).
\end{proof}

The collar lies on a chosen norm box. The closed positive form that
controls a test is selected from that test's finite moment presentation.



\begin{lemma}[Green exchange on a common collar]
\label[lemma]{lem:collar-green}
Let $Q$ be a symmetric positive $d'$- and $d''$-closed current of
bidegree $(n-1,n-1)$ on $\Omega$. Let $\chi\in A_c^{0,0}(\Omega)$.
Suppose that $a$ is smooth, $g$ is locally bounded Borel, and both equal
one smooth function $b$ on a neighborhood of
$\operatorname{supp}d'\chi\cup\operatorname{supp}d''\chi$.
Whenever $D(gQ)$ has coefficient Radon type, one has
\[
 \langle aD(gQ),\chi\rangle
 =\langle gD(aQ),\chi\rangle.
\]
In this statement positivity of $a$ and $g$ is not required.
\end{lemma}

\begin{proof}
Since $Q$ is closed and $a$ is smooth, $D(aQ)=Da\wedge Q$. The
definition of the distributional derivative gives
\begin{align*}
 \langle aD(gQ)-gD(aQ),\chi\rangle
  &=\langle gQ,D(\chi a)-\chi Da\rangle\\
  &=\langle Q,g\{aD\chi+d'\chi\wedge d''a
                         +d'a\wedge d''\chi\}\rangle.
\end{align*}
Every term in braces is supported in the collar. Substituting $a=g=b$
there makes the last expression
\[
 \langle bQ,D(\chi b)-\chi Db\rangle=0.
\]
The final equality follows by applying the same distributional identity
to the smooth pair $(b,b)$. The substitutions are valid identities of
coefficient measures. 
\end{proof}

\begin{lemma}[Monotonicity of the ordered trace integral]
\label[lemma]{lem:trace-monotonicity}
Let $S$ have bidegree $(p,p)$, let $0\leq k\leq n-p$, and let
$\alpha$ be a smooth, positive, $d'$- and $d''$-closed form of bidegree
$(n-p-k,n-p-k)$ such that wedging it with a positive current preserves
positivity. Let $\chi\geq0$ be compactly supported in $\Omega$.

For $0\leq i\leq k$, suppose $f_i,g_i\in\Cb(\Omega)$,
$f_i\leq g_i$, and all these functions equal the same smooth psh
function $b$ near the derivative support of $\chi$. Assume that each
$g_i$ has the smooth majorants with this exact collar supplied by
\cref{lem:smooth-collar}. Define, using the ordered products,
\[
 J(f_0;f_1,\ldots,f_k)
   =\langle f_0T_{S,k}(f_1,\ldots,f_k),\chi\alpha\rangle.
\]
Then
\[
 J(f_0;f_1,\ldots,f_k)
 \leq J(g_0;g_1,\ldots,g_k).
\]

\end{lemma}

\begin{proof}
The case $k=0$ is measure monotonicity. For $k>0$, fix smooth psh
majorants $G_i\geq g_i$ with the common collar. First replace the weight
$f_0$ by $G_0$, using positivity of
$T_{S,k}(f_1,\ldots,f_k)\wedge\alpha$. Apply
\cref{lem:collar-green} with
\[
 Q=T_{S,k-1}(f_2,\ldots,f_k)\wedge\alpha
\]
to move $G_0$ into a differentiated position and $f_1$ into the weight.
The resulting current $DG_0\wedge Q$ is positive and closed. Replace
the weight $f_1$ by $G_1$, and repeat. At the end one obtains
\[
 J(f_0;f_1,\ldots,f_k)
 \leq\langle G_k\,DG_0\wedge\cdots\wedge DG_{k-1}\wedge S,
                     \chi\alpha\rangle.
\]
A final application of the same Green identity, together with
commutation of smooth even-degree factors, makes the right-hand side
$J(G_0;G_1,\ldots,G_k)$.

Now let the index of $G_k$ tend to infinity, keeping every other index
fixed. Bounded convergence against the fixed coefficient measures of
$S$ implies $D(G_kS)\to D(g_kS)$. The outer factors and the weight are
smooth and fixed, so this convergence can be tested directly. Next let
the index of $G_{k-1}$ tend to infinity. The background
$D(g_kS)$ is now fixed, so the same argument applies. Continue in the
order $k,k-1,\ldots,1$. Finally, with the ordered product fixed, let the
weight $G_0$ tend to $g_0$ by bounded convergence. The result is exactly
$J(g_0;g_1,\ldots,g_k)$ in its original order. The order of these limits keeps the coefficient measure fixed at each
discontinuous multiplication.
\end{proof}

\begin{lemma}[An upper bound for decreasing weights]
\label[lemma]{lem:upper-weight}
Let $P_j,P$ be positive coefficient Radon currents of one bidegree,
with $P_j\to P$ on all original compact smooth tests. Let $a_j\downarrow a$
be locally uniformly bounded functions. Suppose that, near the support
of a fixed positive compact test $\Theta$, every fixed $a_l$ has bounded
continuous majorants converging pointwise to $a_l$.
Then
\[
 \limsup_{j\to\infty}\langle a_jP_j,\Theta\rangle
       \leq\langle aP,\Theta\rangle.
\]
\end{lemma}

\begin{proof}
For $j\geq l$ and a continuous majorant $a_{l,\nu}\geq a_l$,
positivity gives
\[
 \langle a_jP_j,\Theta\rangle
       \leq\langle a_{l,\nu}P_j,\Theta\rangle.
\]
Keep $l,\nu$ fixed and pass to the limit in $j$, using \cref{lem:finite-control}. Bounded convergence against the fixed positive Radon measure
$\mu_{P,\Theta}$ first allows $\nu\to\infty$ and then $l\to\infty$.
\end{proof}

\begin{theorem}[Weighted decreasing continuity for a fixed arbitrary background]
\label[theorem]{thm:relative-monotone}
Let $S$ be any
symmetric positive original CLD current of bidegree $(p,p)$ on $U$,
separately $d'$- and $d''$-closed. Let $0\leq k\leq n-p$.
If $u_{i,j}\downarrow u_i$ in $\Cb(U)$ for $0\leq i\leq k$, with
locally bounded limits, then
\[
 u_{0,j}T_{S,k}(u_{1,j},\ldots,u_{k,j})
 \longrightarrow
 u_0T_{S,k}(u_1,\ldots,u_k)
\]
on every original compact smooth complementary test form. The ordered
unweighted products have the corresponding joint decreasing continuity.

\end{theorem}

\begin{proof}
Write $W_k$ for the weighted assertion with $k$ differentiated slots.
Bounded convergence for the fixed coefficient measures of $S$ proves
$W_0$. Distributional differentiation shows that $W_{k-1}$ implies
joint unweighted convergence in degree $k$. We prove $W_k$ assuming
this unweighted convergence.

The statement is local on compact test supports. Use a finite original
smooth partition to reduce to one test $\eta$ supported in a small
norm box $C_1$ as in \cref{lem:smooth-collar}. Include in the chosen
ordinary neighborhood all the unit coordinates needed to present
$\eta$. Let these coordinates be $t_1,\ldots,t_N$, and set
\[
 \beta=\sum_{a=1}^N d't_a\wedge d''t_a,
 \qquad r=n-p-k,\qquad \alpha=\frac{\beta^r}{r!}.
\]
For $r=0$ set $\alpha=1$. The form $\alpha$ is closed and strongly
positive. In the fixed finite coefficient space it is a positive
definite trace form. After replacing $\eta$ by its symmetric part,
there is $C>0$ such that
\[
 \Theta_\pm=C\chi\alpha\pm\eta
\]
are positive original compact tests. This is precisely the matrix
domination used in \cref{lem:finite-control}; additional scalar coordinates
are handled by pullback to a common refinement. 

Patch all the potentials to $v_{i,j},v_i$ by
\cref{lem:smooth-collar}, and put
\[
 P_j=T_{S,k}(v_{1,j},\ldots,v_{k,j}),\qquad
 P=T_{S,k}(v_1,\ldots,v_k).
\]
The induction hypothesis gives $P_j\to P$. Applying
\cref{lem:trace-monotonicity} to the ordered tuples
$(v_0,\ldots,v_k)\leq(v_{0,j},\ldots,v_{k,j})$ yields
\[
 \langle v_0P,\chi\alpha\rangle
       \leq\langle v_{0,j}P_j,\chi\alpha\rangle.
\]
\cref{lem:upper-weight} gives the reverse limsup inequality.
Consequently
\begin{equation}
 \label{eq:trace-limit}
 \langle v_{0,j}P_j,\chi\alpha\rangle
       \longrightarrow\langle v_0P,\chi\alpha\rangle.
\end{equation}
Apply \cref{lem:upper-weight} to $\Theta_+$ and $\Theta_-$.
Together with \eqref{eq:trace-limit}, the two resulting inequalities
give respectively the upper and lower bounds for
$\langle v_{0,j}P_j,\eta\rangle$. Thus this pairing converges to
$\langle v_0P,\eta\rangle$.

On a neighborhood of $\operatorname{supp}\eta$, the patched potentials
are $u_{i,j}+c$ and $u_i+c$. Ordinary locality of the recursion identifies
the differentiated products there with the original products, since
constants in differentiated slots vanish. Remove the constant $c$ from
the weight by the unweighted degree-$k$ convergence already established.
Finite summation over the test partition proves $W_k$.

This induction proves all $W_k$ and all allowed unweighted assertions.

\end{proof}

\begin{corollary}[Symmetry and the first-stage calculus]
\label[corollary]{cor:full-calculus}
Under the hypotheses of \cref{thm:relative-monotone}, the currents
$T_{S,k}$ are symmetric in the differentiated arguments,
positive, symmetric as currents, and separately $d'$- and $d''$-closed.
They are additive and positively homogeneous in each potential, commute
with ordinary restriction and gluing, and are independent of the
regularization used to compute them. In particular,
\[
 D\bigl(uD(vS)\bigr)=D\bigl(vD(uS)\bigr)
 \qquad(u,v\in\Cb(U),\ p\leq n-2).
\]
\end{corollary}

\begin{proof}
All assertions except permutation symmetry were proved for the ordered
recursion. Choose, locally on one fixed domain, continuous semipositive
PL approximants decreasing to each potential, as in \cref{prop:regularization}.
Continuous arbitrary-background products are symmetric, by uniform
smooth approximation and fixed-continuous-multiplier convergence.
\cref{thm:relative-monotone}, with a constant weight, passes this
identity to the given potentials. The recursion identifies the result
independently of the approximants and on overlaps.
\end{proof}

\begin{corollary}[Top-degree convergence as measures]
\label[corollary]{cor:top-measures}
If $p+k=n$, the weighted convergence in
\cref{thm:relative-monotone} holds against every compactly
supported continuous scalar function, locally as signed Radon measures.
\end{corollary}

\begin{proof}
Unweighted convergence against one nonnegative smooth cutoff equal to
one near a prescribed compact set bounds the masses of the positive
measures there. The weights have a common actual pointwise bound on a
slightly larger compact neighborhood. Thus the weighted variations are
uniformly bounded. Uniform density of original smooth scalar functions,
with a fixed cutoff, extends convergence to continuous scalar tests.
\end{proof}


% END sections/03_weighted_continuity.tex

% BEGIN sections/04_curvature_quotient.tex
\section{The curvature quotient of bounded classical potentials}
\label{sec:curvature-quotient}

Write $\DSH(U)$ for the actual real functions locally expressible as
differences of two members of $\Cb$. Let $\PH(U)$ consist of finite
continuous functions whose two signs are classical psh. The following
result from the ordered-potential manuscript is the regularity input
for this section \cite[Theorem 21.4]{Order}.

\begin{theorem}[The bounded-difference criterion]\label[theorem]{thm:difference-criterion}
For $h\in\DSH(U)$,
\[
                  D[h]\ge0\quad\Longleftrightarrow\quad h\in\Cb(U).
\]
Both sides concern the given actual function.
\end{theorem}

We recall the two geometric steps behind the criterion, since its use
is stronger than a choice of a good representative for a distribution.
Order transfer along analytic curve maps after every complete extension
preserves the actual bounded classical difference and the original
curvature inequality. On a curve the positive difference of the two
Laplacians gives an actual subharmonic difference: positive Green
potentials identify its endpoint values along nested discs. To recover
ambient upper semicontinuity, use a compact semistable cover of a norm
box and its finite Shilov interpolants. If $h=a-b$ and $a<a(x)+\epsilon$
on the box, the one-sided curve maximum principle gives
$h\circ q\le P h$ and $b\circ q\le P b$ on each compact chart.
The fiber maximum of the continuous function $P h$ is upper
semicontinuous and majorizes $h$, while its value above $x$ is at most
$h(x)+\epsilon$. Letting $\epsilon$ decrease proves ambient upper
semicontinuity. The full complete-extension curve condition then gives
classical psh membership. These steps and the required algebraic
localization and trace arguments are proved in
\cite[Sections 19--21]{Order}; the zero-curvature predecessor appears
in \cite[Section 10]{Quotient}.

\begin{theorem}[The actual curvature kernel]\label[theorem]{thm:actual-kernel}
The equality
\[
 \ker\bigl(D:\DSH(U)\longrightarrow\mathcal D'^{1,1}(U)\bigr)
                           =\PH(U)
\]
holds for actual functions. In particular, if $u,v\in\Cb(U)$ and
$D[u]=D[v]$, the difference $u-v$ is continuous and both its signs
are classical psh.
\end{theorem}

\begin{proof}
If $D[h]=0$, apply \cref{thm:difference-criterion} to $h$ and $-h$.
Both are upper semicontinuous, so $h$ is continuous; they are classical
psh by the same criterion. Conversely, for $h\in\PH$, first-order
positivity gives $D[h]\ge0$ and $-D[h]\ge0$. Pointedness of the cone
in \cref{lem:finite-control} makes $D[h]=0$. This reasoning is local
and hence applies to every locally represented difference.
\end{proof}

\begin{theorem}[Background invariance and order]
\label[theorem]{thm:background-order}
Let $S$ be a background of degree $(p,p)$, with $p<n$. For
$u,v\in\Cb(U)$,
\[
 D[u]\ge D[v]\quad\Longrightarrow\quad D(uS)\ge D(vS).
\]
Equality of the ambient curvatures gives equality of the displayed
relative curvatures. The same replacement is valid in every
differentiated slot of a mixed product whose other slots are fixed.
\end{theorem}

\begin{proof}
The actual difference $h=u-v$ has positive curvature, hence is classical
psh by \cref{thm:difference-criterion}. The first-order construction
gives $D(hS)\ge0$, and linearity identifies this with
$D(uS)-D(vS)$. For equality, both $h$ and $-h$ give positive relative
curvatures; pointedness makes them zero. To change one mixed slot,
use symmetry to put it first. The product of the remaining slots with
$S$ is one fixed positive separately closed background, to which the
first-order assertion applies.
\end{proof}

This proves the last assertion of \cref{thm:main-bounded}. It also
constructs products of curvature currents having bounded classical
local potentials.

\begin{corollary}[Gluing local-potential currents]
\label[corollary]{cor:potential-current-products}
Suppose each $\alpha_i$ is a $(1,1)$-current locally of the form
$D[a_i]$ with $a_i\in\Cb$. For $k\le n-p$ there is a unique positive
separately closed current
$\alpha_1\wedge\cdots\wedge\alpha_k\wedge S$ whose expression on
such a neighborhood is $T_{S,k}(a_1,\ldots,a_k)$. It is symmetric,
local, additive and positively homogeneous in these current slots.
\end{corollary}

\begin{proof}
On an overlap, corresponding potentials have the same ambient
curvature. Replace the slots one at a time using
\cref{thm:background-order}. The local currents therefore agree on
all original tests and glue by the sheaf property of currents. Their
remaining properties follow locally from the bounded calculus.
\end{proof}

\begin{corollary}[Continuous curvature domination]
\label[corollary]{cor:continuous-domination}
Let $u\in\Cb(U)$, let $g$ be continuous classical psh, and suppose
$0\le D[u]\le C D[g]$ for $C\ge0$. Then $u$ and $Cg-u$ are
continuous classical psh functions.
\end{corollary}

\begin{proof}
The bounded classical difference $Cg-u$ has positive curvature, so it
is upper semicontinuous and classical. The identity
$u=Cg-(Cg-u)$ makes $u$ lower semicontinuous. It is already upper
semicontinuous, hence continuous, and the same follows for $Cg-u$.
\end{proof}

An undifferentiated weight retains its actual value. For example,
$D[0]=D[1]$, while $0S=0$ and $1S=S$. The quotient constructed here
is therefore precisely the quotient in the differentiated slots.

% END sections/04_curvature_quotient.tex

% BEGIN sections/05_continuous_module.tex
\section{Continuous coefficients and mixed curvature operations}
\label{sec:continuous-module}

\begin{theorem}[Transfer of continuous coefficient relations]
\label[theorem]{thm:continuous-relations}
Let $g_1,\ldots,g_m\in\Cb(U)$ and $f_1,\ldots,f_m\in C(U,\Real)$.
For a fixed background $S$, put
\[
 A=\sum_i f_iD[g_i],\qquad A_S=\sum_i f_iD(g_iS).
\]
Then $A\ge0$ implies $A_S\ge0$. In particular, $A=0$ implies
$A_S=0$ on every original complementary test.
\end{theorem}

\begin{proof}
Set $T=\sum_iD(g_iS)\ge0$. For $\epsilon>0$, work near a point on
a neighborhood where $|f_i-c_i|<\epsilon$ for suitable real constants.
There,
\[
 D\left[\sum_i(c_i+\epsilon)g_i\right]
       =A+\sum_i(c_i+\epsilon-f_i)D[g_i]\ge0.
\]
The actual function inside brackets is a bounded classical difference,
as is seen by separating its positive and negative real coefficients.
By \cref{thm:difference-criterion} it is classical psh. Its curvature
against $S$ is positive, so $\sum_i c_iD(g_iS)\ge-\epsilon T$.
The coefficient error gives $A_S\ge-2\epsilon T$ on this neighborhood.

Cover the support of a positive original test by finitely many such
neighborhoods and use a nonnegative smooth partition. Summing gives the
same inequality on that test. Its control mass against $T$ is finite,
so $\epsilon\downarrow0$ proves $A_S\ge0$. If $A=0$, apply the
argument to both coefficient lists $(f_i)$ and $(-f_i)$ and use
pointedness of the positive cone.
\end{proof}

We now prove the finite-presentation part of \cref{thm:main-module}.

\begin{proof}[Proof of the continuous coefficient operation]
Each expression in \MathEqref{eq:main-module-action} is a finite sum
of continuous multiples of positive coefficient Radon currents. If a
presentation of one slot changes without changing its actual current,
fix a multi-index in the other slots and set
\[
 R=T_{S,r-1}(g_{1,\alpha_1},\ldots,
                \widehat{g_{i,\alpha_i}},\ldots,g_{r,\alpha_r}).
\]
This is a fixed positive separately closed background. The zero-relation
case of \cref{thm:continuous-relations} makes the change in that slot
zero. Multiply by the remaining continuous coefficients and sum. Changing
one slot at a time proves independence of all presentations. Symmetry,
multilinearity and ordinary locality follow from the bounded products.

For positivity, fix a positive compact complementary test and localize
to one finite presentation. Its coefficients are bounded by $M\ge1$
on a slightly larger compact set. Subdivide the test support into
ordinary neighborhoods on which
$|f_{i,\alpha}-c_{i,\alpha}|<\epsilon$, with $0<\epsilon<1$.
On one such neighborhood put
\[
 H_i=\sum_\alpha(c_{i,\alpha}+\epsilon)g_{i,\alpha}.
\]
If $A_i\ge0$, then
\[
 D[H_i]=A_i+\sum_\alpha(c_{i,\alpha}+\epsilon-f_{i,\alpha})
                                             D[g_{i,\alpha}]\ge0.
\]
Thus $H_i\in\Cb$ by \cref{thm:difference-criterion}, and
$T_{S,r}(H_1,\ldots,H_r)\ge0$.

For each multi-index, telescope the product of its $r$ scalar
coefficients. The difference between the coefficient in
\MathEqref{eq:main-module-action} and that in the last positive
product is at most $2r\epsilon(M+2)^{r-1}$ in absolute value.
Let $R_0$ be the finite sum of all the positive BT tails in that
presentation, with their scalar coefficients omitted. Hence
\[
       \Phi_{S,r}(A_1,\ldots,A_r)
            \ge-2r\epsilon(M+2)^{r-1}R_0
\]
on each member of the subdivision. A nonnegative partition and the
fixed finite control mass pass this inequality to the original test.
Letting $\epsilon$ decrease proves positivity. Telescoping in the slots
then proves monotonicity when the remaining slots are positive.
\end{proof}

\begin{example}[A Borel coefficient can destroy a relation]
\label[example]{ex:borel-relation}
On the bidannulus
$U=\{r<|T_1|<R,\ r<|T_2|<R\}$, where $0<r<1<R$, take
$Y=\{T_2=1\}\cap U$ and $g=\max(0,\log|T_1|)$.
The first-order analytic-divisor nonconcentration theorem gives
$\mathbf1_YD[g]=0$ as an ambient coefficient Radon current.
The subvariety formula and the slope jump on the $T_1$-annulus give
\[
 \mathbf1_YD(g[Y])=D(g[Y])=i_*\delta_{\zeta_{0,1}}\ne0,
 \qquad i:Y\hookrightarrow U.
\]
Here $\zeta_{0,1}$ is the radius-one Gauss point; the slope jump is one.
A compact nonnegative original scalar test equal to one near that point
has integral one. This is the obstruction to Borel coefficient transfer described in
\cite[Section \emph{Products in the continuous coefficient curvature module}]{Order}.
\end{example}

Continuous variable coefficients also affect closedness. The operation
is an operation on coefficient currents; an output serves as a new
background when its separate closedness has been established.

% END sections/05_continuous_module.tex

% BEGIN sections/06_order_completion.tex
\section{Order completion inside the original currents}
\label{sec:order-completion}

The continuous coefficient operation has a useful continuity stronger
than distributional continuity. A fixed positive control measures the
size of a signed current and, simultaneously, the error of all its
background actions. We follow the intrinsic order construction of
\cite{Intrinsic}.

\begin{lemma}[Signed products]\label[lemma]{lem:signed-order}
Let $Q_i\in\Mone$ be positive and suppose
$-M_iQ_i\le A_i\le M_iQ_i$, where $M_i\ge0$. Then
\begin{equation}\label{eq:signed-order}
 -\Bigl(\prod_iM_i\Bigr)R_S
 \le\Phi_{S,r}(A_1,\ldots,A_r)
 \le\Bigl(\prod_iM_i\Bigr)R_S,
 \qquad R_S=\Phi_{S,r}(Q_1,\ldots,Q_r).
\end{equation}
The difference estimate \MathEqref{eq:main-order-estimate} also holds.
\end{lemma}

\begin{proof}
Set $A_i^\pm=(M_iQ_i\pm A_i)/2$. These are positive,
$A_i=A_i^+-A_i^-$, and $A_i^++A_i^-=M_iQ_i$. Expand the mixed
product into a signed sum of positive currents. The unsigned sum is
$\Phi_{S,r}(M_1Q_1,\ldots,M_rQ_r)$. This unsigned sum plus or minus
the signed sum is twice a positive subsum, proving
\MathEqref{eq:signed-order}. Telescope the difference between two
mixed products one slot at a time and apply this inequality to each
term. Addition gives \MathEqref{eq:main-order-estimate}.
\end{proof}

For one positive $Q\in\Mone(U)$, the space $E_Q$ and its gauge were
defined in the introduction. Every finite presentation belongs locally
to some such space: bound its coefficients and sum its positive
curvature generators.

\begin{theorem}[Faithful realization]\label[theorem]{thm:faithful-order}
The completion $\widehat E_Q(U)$ embeds isometrically into the actual
symmetric coefficient Radon $(1,1)$-currents. Its image is the closure
of $E_Q(U)$ for the ambient order norm defined by $Q$. If an actual
current is obtained using controls $Q$ and $R$, its two realizations
agree in the completion controlled by $Q+R$.
\end{theorem}

\begin{proof}
If $\|A\|_Q=0$, then $-\epsilon Q\le A\le\epsilon Q$ for every
$\epsilon>0$. The finite test control makes every pairing of $A$
zero. Thus the gauge is a norm. Addition of order bounds gives its
triangle inequality. Weak closedness of the positive cone shows that
the infimum in a finite gauge is attained.

Let $(A_j)$ be order Cauchy. On each fixed original compact test its
values are Cauchy by \cref{lem:finite-control}. The finite coefficient
bound makes their limit a distribution and a coefficient Radon current
on every compact presentation. Symmetry and compatibility with all
relations and refinements pass to this limit, denoted $A$. If
$\|A_j-A_l\|_Q\le\epsilon$ for $j,l\ge N$, let $l$ tend to infinity
in both order inequalities. Weak closedness gives
\begin{equation}\label{eq:order-limit}
                -\epsilon Q\le A_j-A\le\epsilon Q\qquad(j\ge N).
\end{equation}
The convergence is therefore convergence in the ambient order norm.
If $A=0$, the same inequality shows that the original sequence
represents zero in the completion. This proves injectivity, the
isometry, and the description of the image.

Since $Q,R\le Q+R$, both completion sequences are Cauchy for that
larger control. Their actual limits agree; injectivity for $Q+R$
identifies the two completion classes.
\end{proof}

An actual symmetric coefficient current belongs to $\Mone_{\rm ord}$
when it belongs locally to one of these embedded completions. This is
a sheaf by the local definition and the sheaf property of currents.
Local addition uses $Q+R$, and multiplication by a continuous function
uses its bound on a smaller neighborhood.

\begin{theorem}[Extension of the mixed operation]
\label[theorem]{thm:extended-operation}
The maps $\Phi_{S,r}$ extend uniquely to symmetric locally
$C(U,\Real)$-multilinear operations on $\Mone_{\rm ord}$, continuous
for every common finite family of local order controls. They are
positive, local, compatible with gluing, and satisfy the signed and
difference estimates. Their supports satisfy
\begin{equation}\label{eq:product-support}
 \supp\Phi_{S,r}(A_1,\ldots,A_r)
       \subset\supp S\cap\bigcap_i\supp A_i.
\end{equation}
If the mixed products along an order-convergent approximation are
separately closed, their limit is separately closed.
\end{theorem}

\begin{proof}
Choose local controls $Q_i$ and Cauchy approximants $A_{i,j}$. Their
order norms are bounded. The difference estimate and finite test
control make the mixed products Cauchy in coefficient measure on every
fixed compact presentation. Their limit is therefore a coefficient
Radon current. The same estimate applied to two approximating families
proves independence. For different controls $R_i$, pass to $Q_i+R_i$
and use \cref{thm:faithful-order}; the resulting limits agree.

Algebraic identities pass to the limit. If $A_i\ge0$, the inequality
\MathEqref{eq:order-limit} gives positive approximants
$A_{i,j}+\epsilon_{i,j}Q_i$, with $\epsilon_{i,j}\to0$. These have the
same limits, so positivity passes from the finite-presentation
operation. The order estimates also pass by weak closedness.

Restriction is contractive for the restricted controls, so these local
constructions agree on intersections and glue. On an open where one
slot is zero, its completion class is zero by faithful realization,
and multilinearity makes the product zero. The same holds where
$S=0$, proving \MathEqref{eq:product-support}. Finally, distributional
continuity of $d'$ and $d''$ proves the closedness assertion.
\end{proof}

This completes \cref{thm:main-module}. The numerical coefficient in
the order estimate is independent of $S$; the controlling trace mass
is evaluated against the particular background and test.

\begin{corollary}[Fixed-generator approximation by curvatures]
\label[corollary]{cor:fixed-generator-closedness}
Suppose locally
\[
 A=\sum_{i=1}^N f_iD[g_i],\qquad
 A_j=\sum_{i=1}^N f_{i,j}D[g_i]=D[v_j],
\]
where $g_i$ and $v_j$ are finite continuous local classical psh
differences, $f_i,f_{i,j}$ are continuous, and $f_{i,j}\to f_i$ locally
uniformly. Then $\Phi_{S,1}(A)$ is separately closed for every allowed
background $S$.
\end{corollary}

\begin{proof}
Write $g_i=g_i^+-g_i^-$ and choose the fixed control
$Q=\sum_i(D[g_i^+]+D[g_i^-])$. On a smaller neighborhood the error
is bounded between $\pm\epsilon_jQ$, where
$\epsilon_j=\max_i\|f_i-f_{i,j}\|_\infty\to0$. Thus the actions
converge on all original tests. Continuous relation transfer identifies
$\Phi_{S,1}(A_j)=D(v_jS)$, a separately closed current. Both first
derivatives pass to the limit. This argument uses the coefficient
convergence and requires no bound or normalization on $v_j$.
\end{proof}

\begin{proposition}[A support-conormal certificate]
\label[proposition]{prop:support-conormal}
Let $A_i=D[g_i]$ for continuous classical psh differences, and set
$T_i=D(g_iS)$ and $F_i=\supp A_i$. Then $\supp T_i\subset F_i$.
If a smooth function $h$ vanishes on $F_i$, then
$d'h\wedge T_i=d''h\wedge T_i=0$.

Suppose smooth covectors admit an identity
\[
 \theta_i=\sum_a f_{ai}\eta_a+\sum_j b_{ij}d'h_{ij},\qquad
 \sum_i f_{ai}A_i=0,\qquad h_{ij}|_{F_i}=0,
\]
with continuous scalar coefficients and original smooth
$\eta_a,h_{ij}$. Then $\sum_i\theta_i\wedge T_i=0$. The same
assertion holds after interchanging the two differential columns.
\end{proposition}

\begin{proof}
On the complement of $F_i$, the zero curvature relation transfers to
$T_i=0$. This proves the support assertion. Hence $hT_i=0$; separate
closedness and the smooth product rule give the conormal identities.
Each scalar relation transfers to $\sum_i f_{ai}T_i=0$. Multiply by
$\eta_a$, apply the conormal identity to the remaining terms, and sum.
\end{proof}

% END sections/06_order_completion.tex

% BEGIN sections/07_relative_energy.tex
\section{Relative power energies on Laurent domains}
\label{sec:relative-energy}

\subsection{Truncation, capacity and the boundary class}

We now impose the strict Laurent-domain and background-extension
hypotheses \MathEqref{eq:main-box}. All capacities use the entire
normalized classical tester class on the indicated domain. Its ordinary
open outer regularization is denoted $C_S^*$. Local truncated capacity
convergence means that, for every ordinary $\Omega\subset V$, compact
$L\Subset\Omega$, $R>0$ and $\delta>0$,
\[
 C_{S|\Omega}^*\bigl(L\cap
 \{|\max(u_j,-R)-\max(u,-R)|>\delta\};\Omega\bigr)\to0.
\]

Let $u_1,\ldots,u_m\le0$ be extended classical potentials, with their
actual values. Put $u_i^{(j)}=\max(u_i,-j)$ and
$E_j=\bigcap_i\{u_i>-j\}$. The restrictions
\begin{equation}\label{eq:nonpolar-cut}
 \mathbf1_{E_j}T_{S,k}(u_1^{(j)},\ldots,u_k^{(j)})
\end{equation}
are consistent on their common finite-value regions. Their increasing
positive coefficient limits define the nonpolar product whenever those
limits are locally finite. For a single potential write
$\mu_{u,j}=T_{S,q}((u^{(j)})^{[q]})$,
$\nu_{u,j}=\mathbf1_{\{u\le-j\}}\mu_{u,j}$, and $\mu_u$ for the
nonpolar top measure.

\begin{definition}[The actual boundary power classes]
\label[definition]{def:power-class}
The boundary condition $\mathrm{BG}(u,0)$ means
$\overline{\{u<-t\}}\Subset V$ for every $t>0$. The class
$\mathcal F_{S,\partial}(V)$ consists of negative extended classical
potentials satisfying this condition and
$\nu_{u,j}(V)\to0$. For $\sigma>0$, set
\[
 \mathcal E^\sigma_{S,\partial}(V)
 =\left\{u\in\mathcal F_{S,\partial}(V):
          e_\sigma(u)=\int_V(-u)^\sigma\,d\mu_u<\infty\right\}.
\]
The integral is over the finite-value locus, to which $\mu_u$ is
carried. Write $\mathcal E^1=\mathcal E^1_{S,\partial}(V)$.
\end{definition}

The following analytic results will be used with their displayed
domain. Their proofs in \cite{Energy} supply the coefficient cutoffs,
capacity comparison, and tail bounds needed for the later completions.

\begin{theorem}[Analytic truncation estimates]\label[theorem]{thm:analytic-estimates}
The following statements hold.
\begin{enumerate}
\item A locally coefficient-finite product
\MathEqref{eq:nonpolar-cut} is separately closed on all original tests.
This is an ordinary-local assertion. Under $\mathrm{BG}(u,0)$,
$\mu_u$ is locally finite and
\[
 u\in\mathcal F_{S,\partial}(V)
 \quad\Longleftrightarrow\quad
 t^qC_S(\{u<-t\};V)\to0.
\]
\item For $\sigma>0$ and $u$ satisfying $\mathrm{BG}(u,0)$, put
\[
 J_\sigma(u)=\int_0^\infty
               t^{q+\sigma-1}C_S(\{u<-t\};V)\,dt.
\]
Then $J_\sigma(u)<\infty$ exactly when
$u\in\mathcal E^\sigma_{S,\partial}(V)$, and
\begin{equation}\label{eq:capacity-energy}
 \frac{\sigma}{2^{q+\sigma}}J_\sigma(u)
       \le e_\sigma(u)
       \le\frac{\sigma2^q}{1-2^{-\sigma}}J_\sigma(u).
\end{equation}
Raw and outer capacities agree on these ordinary open sublevels.
\item For a finite tuple $\mathbf u$ in this power class and a
multi-index $\gamma$ with $k=|\gamma|\le q$, its unbounded mixed
curvature $N_\gamma$ is defined, positive and separately closed.
For $k\le q-1$,
\begin{equation}\label{eq:actual-recursion}
 D(u_iN_\gamma)=N_{\gamma+e_i},\qquad
 \mathcal J'_{\gamma,i}=d'(u_iN_\gamma),\quad
 \mathcal J''_{\gamma,i}=d''(u_iN_\gamma).
\end{equation}
The first-order currents are of coefficient Radon type. When
$k+2\le q+\sigma$, the energy matrix is also of coefficient Radon
type and is given by
\begin{equation}\label{eq:three-term-energy}
 G^\gamma_{ih}
   =\tfrac12\left[D(u_i u_hN_\gamma)
       -u_iN_{\gamma+e_h}-u_hN_{\gamma+e_i}\right].
\end{equation}
It is positive semidefinite in the middle cone.
\item Put $\kappa=q+\sigma$, $w=\sum_i u_i$,
$W=1-w$, and
$J_w(R)=\int_R^\infty t^{\kappa-1}C_S(\{w<-t\};V)\,dt$.
For a curvature, first-order or energy coefficient current $P$, assign
the degree budget $b=k,k+1,k+2$, respectively. If $a+b\le\kappa$,
the actual and canonically truncated coefficients have the uniform
tail bound
\begin{equation}\label{eq:graded-tail}
 \int_{\{w\le-R\}}W^a\,|d\mu_{P,\eta}|
       \le C_\eta R^{a+b-\kappa}J_w(R/4),\qquad R\ge1,
\end{equation}
with the corresponding bound for every bounded truncation.
Canonical truncations converge in local coefficient total variation,
including continuous tuple multipliers of growth at most $W^a$.
\item Bounded actual Borel capacity limits of classical functions in
$[-R,0]$ have an intrinsic bounded calculus, including first-order and
energy currents. Its degree-$b$ coefficients satisfy
\begin{equation}\label{eq:bounded-capacity-control}
                |\mu_{P_R,\eta}|(E)\le C_\eta R^bC_S^*(E;V).
\end{equation}
These currents are independent of defining rows and outer-null changes,
and bounded continuous tuple weights pass under capacity convergence.
\item There is a bounded continuous classical barrier
\begin{equation}\label{eq:box-barrier}
 \rho=\max\left(-1,\max_a\log\frac{|f_a|}{R_a},
                      \max_b\log\frac{r_b}{|g_b|}\right)
\end{equation}
on a neighborhood of $K_0$, with $-1\le\rho<0$ on $V$, compact
negative sublevels, and $\mu_\rho(V)<\infty$. For
$u\in\mathcal E^\sigma$, the bounded finite-mass potentials
$\max(u,j\rho)$ decrease to $u$. Fixed lower-barrier families have
critical weighted convergence, including global mixed-energy
convergence in first energy.
\end{enumerate}
\end{theorem}

The individual restricted stages in \MathEqref{eq:nonpolar-cut} may
have boundary derivatives; the closedness theorem concerns their
locally finite limit. The exact sources of
\cref{thm:analytic-estimates} are
\cite[Sections 8, 12, 15--20, 31]{Energy}. In particular, the
first-order estimate is (FO8), the curvature and energy bounds are
(PG3) and (PG12), and the bounded Borel construction is the lemma
\emph{Analytic bounded capacitary calculus}. The boundary comparison
and all-degree capacity estimates also give the following consequences
used below: equality of two no-loss top measures implies equality of
their potentials modulo local outer capacity; the actual power classes
are closed under finite sums and maxima; and capacity convergence of
uniformly energy-bounded tuples passes bounded weighted products below
the critical degree. These are proved in
\cite[Sections 12, 14--18, 28]{Energy}.

\subsection{Global Green symmetry and the mixed energy inequality}

For bounded classical $a,b$ and a closed background $Q$ of degree
$(n-1,n-1)$, write
$\Gamma_Q(a,b)=\frac12[D(abQ)-aD(bQ)-bD(aQ)]$.
The bounded energy identity and its positive energy matrix are those
of the relative first-order calculus \cite{Relative}; the same
definition is justified for the unbounded products by
\MathEqref{eq:three-term-energy}.

\begin{lemma}[Global bounded Green exchange]\label[lemma]{lem:global-green}
Let $a,b\le0$ be bounded classical potentials with compact negative
sublevel closures. If $D(aQ)$ and $D(bQ)$ have finite total mass, then
\[
 \int_V(-a)D(bQ)=\int_V(-b)D(aQ)=\int_V\Gamma_Q(a,b).
\]
The diagonal energy is finite, and the cross energy has finite total
variation.
\end{lemma}

\begin{proof}
Put $a^\epsilon=\max(a,-\epsilon)$. The difference
$a-a^\epsilon$ is supported in
$K_\epsilon=\overline{\{a<-\epsilon\}}\Subset V$. Compact Green
exchange gives
\[
 \int(a-a^\epsilon)D(bQ)
       =\int bD(aQ)-\int bD(a^\epsilon Q).
\]
The measures $D(a^\epsilon Q)$ and $D(aQ)$ have the same finite
total mass: their difference is the derivative of a compactly
supported current, which pairs to zero with a cutoff equal to one
near that support. For fixed $\delta>0$, relative CLN on
$L_\delta=\overline{\{b<-\delta\}}$ gives
$D(a^\epsilon Q)(L_\delta)\le C_{\delta,Q}\epsilon$. Therefore
\[
 \left|\int bD(a^\epsilon Q)\right|
       \le\|b\|_\infty C_{\delta,Q}\epsilon
                               +\delta D(aQ)(V).
\]
Let $\epsilon\downarrow0$ and then $\delta\downarrow0$ to obtain
the first equality.

The bounded clipping rule and the zero energy of an exact scalar
level give
\[
 \int_{\{a<-\epsilon\}}\Gamma_Q(a)
                     =\int(a^\epsilon-a)D(aQ).
\]
Monotone convergence identifies its limit with
$\int(-a)D(aQ)<\infty$. Apply the diagonal identity also to $b$ and
$a+b$ and polarize. Measure Cauchy--Schwarz bounds the total variation
of the cross energy by the square root of the two diagonal masses.
\end{proof}

\begin{theorem}[First mixed energy]\label[theorem]{thm:first-mixed-energy}
For $u_0,\ldots,u_q\in\mathcal E^1$, the finite functional
\[
 B(u_0,\ldots,u_q)=\int_V(-u_0)N_q(u_1,\ldots,u_q)
\]
is symmetric and separately additive and positively homogeneous.
For $Q=N_{q-1}(c_2,\ldots,c_q)$,
\begin{equation}\label{eq:global-energy-matrix}
 \int_V\Gamma_Q(a,b)=B(a,b,c_2,\ldots,c_q),\qquad
 |B(a,b,\mathbf c)|^2\le B(a,a,\mathbf c)B(b,b,\mathbf c).
\end{equation}
The energy matrix stays positive on finite real linear combinations
of the potentials, and all cross energies have finite total variation.
Moreover,
\begin{equation}\label{eq:holder}
 B(u_0,\ldots,u_q)\le\prod_{i=0}^q e(u_i)^{1/d},\qquad
 e(u+v)^{1/d}\le e(u)^{1/d}+e(v)^{1/d}.
\end{equation}
The constant one is optimal whenever the class contains a positive
energy potential.
\end{theorem}

\begin{proof}
\textbf{Step 1: exchange the undifferentiated slot.}
Use $b_{i,j}=\max(u_i,j\rho)$. For fixed $j$ these potentials are
bounded and their sum is bounded below by $(q+1)j\rho$. Boundary
comparison with $c(q+1)j\rho$, followed by $c\downarrow1$, gives
\[
 \mu_{\sum_i b_{i,j}}(V)
       \le[(q+1)j]^q\mu_\rho(V)<\infty.
\]
The positive mixed expansion makes each top mixed measure finite.
Apply \cref{lem:global-green} after fixing the last $q-1$ slots.
Together with differentiated-slot symmetry this proves full bounded
mixed symmetry. The fixed lower-barrier convergence in
\cref{thm:analytic-estimates} passes the identity globally to the
$u_i$. Finiteness is the mixed power-energy estimate in
\cite[(PE11)]{Energy}. Additivity follows from integration and the
unbounded recursion.

\textbf{Step 2: identify the global energy measures.}
For $a^\epsilon=\max(a,-\epsilon)$, apply compact Green to the
canonical bounded difference $a^{(j)}-a^\epsilon$ and to the bounded
background $Q_j$. The clipping rule gives
\[
 \int_{\{a<-\epsilon\}}\Gamma_{Q_j}(a^{(j)})
       =\int(a^\epsilon-a^{(j)})D(a^{(j)}Q_j).
\]
The support lies in a fixed compact set. Critical local total-variation
convergence from \cref{thm:analytic-estimates} passes the energy side;
the degree-one weighted curvature convergence passes the other side.
Thus the same identity holds for $a,Q$. Let $\epsilon\downarrow0$.
The energy is zero on the pole and the scalar level $a=0$, and the
right side tends to its finite mixed energy. This proves the diagonal
identity. Polarization and positive energy matrices give
\MathEqref{eq:global-energy-matrix} and the variation bound.

\textbf{Step 3: prove the sharp integral inequality.}
The zero-factor case follows from the mixed power-energy estimate.
Otherwise normalize $e(u_i)=1$. Among the finitely many mixed values
$B_\alpha=B(u_0^{[\alpha_0]},\ldots,u_q^{[\alpha_q]})$,
$\sum\alpha_i=d$, choose a largest value $M$, and among its indices
maximize $\sum\alpha_i^2$. If two entries $\alpha_i,\alpha_l$ are
positive, energy Cauchy--Schwarz gives
\[
 M^2\le B_{\alpha+e_i-e_l}B_{\alpha-e_i+e_l}\le M^2.
\]
Since a diagonal value is one, $M>0$, and both neighboring values
equal $M$. At least one neighboring index has larger sum of squares,
a contradiction. Thus a maximizing index is pure and $M=1$.
Rescaling proves the first inequality in \MathEqref{eq:holder}.
Expand $e(u+v)$ and use the binomial theorem for the second. Equality
with all slots equal proves optimality.
\end{proof}

\subsection{Discrepancy and quantitative tests}

Set $\mathscr A(u)=-e(u)/d$. Symmetric multilinearity and global
Green give, for $h=v-u$ and $u_t=(1-t)u+tv$,
\begin{equation}\label{eq:energy-variation}
 \frac d{dt}\mathscr A(u_t)=\int h\,d\mu_{u_t},\qquad
 \frac{d^2}{dt^2}\mathscr A(u_t)
        =-q\int\Gamma_{N(u_t^{[q-1]})}(h)\le0.
\end{equation}
In particular
\begin{equation}\label{eq:discrepancy}
 I(u,v)=e(u)+e(v)-B(u,v^{[q]})-B(v,u^{[q]})
       =\sum_{l=0}^{q-1}\int\Gamma_{N(u^{[l]},v^{[q-1-l]})}(u-v).
\end{equation}
It is nonnegative and vanishes precisely on equal relative capacity
classes. To see the latter assertion, zero discrepancy makes each
diagonal energy in the sum zero. Compact Green and Cauchy--Schwarz,
passed through canonical truncations, then give
$D((u-v)N(u^{[l]},v^{[q-1-l]}))=0$ on every test. Summing the
telescoping identity gives $\mu_u=\mu_v$. The boundary comparison
consequence of \cref{thm:analytic-estimates} identifies their relative
classes. The reverse follows from representative independence.

\begin{theorem}[Uniform energy tests]\label[theorem]{thm:uniform-energy-test}
Put $\alpha=2^{-q}$. For $u,v,\psi\in\mathcal E^1$ and
$M=\max\{e(u+v),e(\psi)\}$,
\begin{equation}\label{eq:uniform-energy-test}
 \int_V|u-v|\,d\mu_\psi
       \le2^{(q-1)\alpha}I(u,v)^\alpha M^{1-\alpha}.
\end{equation}
For a mixed top measure, replace $\psi$ by the sum of its $q$
potentials and divide the bound by $q!$. At $M=0$ the integral is zero.
\end{theorem}

\begin{proof}
Let $h=u-v$, $w=u+v$, and
$J_k=\int\Gamma_{N(\psi^{[k]},w^{[q-1-k]})}(h)$ for
$0\le k\le q-1$. Positive expansion gives
$J_0\le2^{q-1}I(u,v)$. Exchange one $\psi$ slot with one $h$ slot by
global symmetry and expand the latter as $u-v$. Cauchy--Schwarz in
the two positive backgrounds, then in their two scalar terms, gives
\[
 J_{k+1}\le
 \{J_k B(\psi^{[k+2]},w^{[q-1-k]})\}^{1/2}
       \le(J_kM)^{1/2}.
\]
Hence $J_{q-1}\le J_0^{2^{-(q-1)}}M^{1-2^{-(q-1)}}$.
The bounded absolute-value chain, the zero energy of an exact scalar
level, and critical total-variation passage give
$\Gamma_Q(|h|)=\Gamma_Q(h)$ in these backgrounds. Global Green gives
\[
 \int|h|\,d\mu_\psi
    =-\int\Gamma_{N(\psi^{[q-1]})}(|h|,\psi)
       \le J_{q-1}^{1/2}e(\psi)^{1/2}.
\]
Combine the inequalities. The zero case follows from mixed energy
finiteness and \MathEqref{eq:holder}. Positive expansion
$q!N(\psi_1,\ldots,\psi_q)\le\mu_{\sum\psi_i}$ proves the last
assertion.
\end{proof}

We record the remaining quantitative consequences needed for the
completion. They are proved in
\cite[Section \emph{Quantitative first-energy convergence in every local capacity}]{Energy}.
For an energy bound $e(u),e(v)\le M_0$ and every compact/domain pair
$L\Subset\Omega\subset V$,
\begin{align}
 C_{S|\Omega}^*(L\cap\{|u-v|>\delta\};\Omega)
       &\le C_{L,\Omega,M_0,S}\delta^{-1}I(u,v)^\alpha,
                  \label{eq:energy-capacity}\\
 |e(u)^{1/d}-e(v)^{1/d}|^d&\le I(u,v),
                  \label{eq:energy-root}\\
 I(u,u^{(R)})&\le\int(u^{(R)}-u)\,d\mu_u\longrightarrow0,
                  \label{eq:clipping-strong}\\
 \int_R^\infty t^qC_S(\{u<-t\};V)\,dt
       &\le2^d\int(-u-R/2)^+\,d\mu_u.
                  \label{eq:capacity-tail-energy}
\end{align}
The capacity estimate uses a normalized tester $a$ replaced near $L$
by $\max(a-2,A_L\rho)$. Its energy is uniformly bounded by the finite
mass of $\rho$, so \MathEqref{eq:uniform-energy-test} applies. Outer
quasicontinuity and the analytic domain comparison give all local outer
capacities. For \MathEqref{eq:energy-root}, write
$a=e(u)^{1/d}$, $b=e(v)^{1/d}$ and use mixed H\"older to bound
$I(u,v)$ below by $(a-b)(a^{d-1}-b^{d-1})\ge|a-b|^d$.
The clipping bound follows from the definition of $I$ and positivity;
dominated convergence passes its right side. The tail estimate is the
tail form of \MathEqref{eq:capacity-energy}.

\begin{proposition}[Sharp midpoint bound]\label[proposition]{prop:midpoint}
For $u,v\in\mathcal E^1$,
\begin{equation}\label{eq:midpoint}
 \mathscr A\left(\frac{u+v}{2}\right)
       -\frac{\mathscr A(u)+\mathscr A(v)}2
       \ge c_q I(u,v),\qquad
 c_q=\frac{1-2^{-q}}{2(q+1)}.
\end{equation}
This constant is optimal if a positive-energy potential exists.
\end{proposition}

\begin{proof}
The scalar Green formula on $[0,1]$ and
\MathEqref{eq:energy-variation} express the left side as
$\sum_{l=0}^{q-1}a_{q,l}\gamma_l$, where
\[
 \gamma_l=\int\Gamma_{N(u^{[l]},v^{[q-1-l]})}(u-v)\ge0,
 \quad
 a_{q,l}=\frac q2\binom{q-1}{l}
           \int_0^1\min(t,1-t)t^l(1-t)^{q-1-l}\,dt.
\]
Their sum $\sum_l\gamma_l$ is $I(u,v)$. We verify
$a_{q,l}\ge c_q$. Put
$p_i(s)=\binom qi s^i(1-s)^{q-i}$,
$B_k(s)=\sum_{i=k}^qp_i(s)$, and
$b_k(t)=q\binom{q-1}{k-1}t^{k-1}(1-t)^{q-k}$.
Differentiation gives $B_k'=b_k$. For $0\le s\le1/2$,
\[
 P_k(s):=\int_s^{1-s}b_k(t)\,dt
       =1-B_k(s)-B_{q-k+1}(s).
\]
If $k\le(q+1)/2$, then
$P_k(s)-P_1(s)=\sum_{i=1}^{k-1}(p_i(s)-p_{q-i}(s))\ge0$,
since $(s/(1-s))^{2i-q}\ge1$; the remaining cases follow from
$P_k=P_{q+1-k}$. Layer cake for the tent function gives
$a_{q,k-1}=\frac12\int_0^{1/2}P_k(s)\,ds\ge a_{q,0}$.
The endpoint integral is $a_{q,0}=c_q$, proving the bound.
For $v=0$ and $e(u)>0$, homogeneity gives $I(u,0)=e(u)$ and
equality in \MathEqref{eq:midpoint}.
\end{proof}

% END sections/07_relative_energy.tex

% BEGIN sections/08_strong_completion.tex
\section{The complete relative first-energy space}
\label{sec:strong-completion}

Retain the domain and background of \MathEqref{eq:main-box}.
Put $d=q+1$, $\alpha=2^{-q}$ and
$\mathcal E^1=\mathcal E^1_{S,\partial}(V)$. We construct
actual Borel classes from the strong Cauchy condition and then define
their currents through bounded truncations.

\subsection{Uniform tails and the boundary of an energy-Cauchy row}

Call an actual row $(u_j)\subset\mathcal E^1$ strongly Cauchy if
\begin{equation}\label{eq:strong-cauchy}
               \lim_{N\to\infty}\sup_{j,l\geq N}
                    I(u_j,u_l)=0.
\end{equation}
The Cauchy condition is stated directly in terms of this discrepancy.

\begin{lemma}[Uniform deep tails and global weighted tightness]
\label[lemma]{lem:r11-strong-cauchy-uniform-tails}
Every strongly Cauchy actual row has bounded first energies and
\begin{equation}\label{eq:uniform-deep-tails}
 \begin{split}
 \lim_{R\to\infty}\sup_j
     \int_{\{u_j<-R\}}(-u_j)\,d\mu_{u_j}&=0,\\
 \lim_{R\to\infty}\sup_j
     \int_R^\infty t^q C_S(\{u_j<-t\};V)\,dt&=0.
 \end{split}
\end{equation}
For finitely many strongly Cauchy actual rows $u_{i,j}$,
$0\leq i\leq q$, the finite positive Radon measures
\begin{equation}\label{eq:cauchy-weighted-measures}
            (-u_{0,j})M(u_{1,j},\ldots,u_{q,j})
\end{equation}
are uniformly tight on $V$. The same conclusions hold if the
indices of the finitely many rows tend independently to infinity.
\end{lemma}

\begin{proof}
The energy-root estimate \MathEqref{eq:energy-root} bounds the energies by some $M$.
The case $M=0$ is immediate from the constant-one mixed H\"older
inequality; otherwise enlarge $M$ to be positive. On
$\{-u_j>R\}$ split according to whether
$-u_N\leq(-u_j)/2$. The pointwise splitting used below gives
\begin{equation}\label{eq:tail-splitting}
 \int_{\{u_j<-R\}}(-u_j)\,d\mu_{u_j}
 \leq 2\int|u_j-u_N|\,d\mu_{u_j}
       +4\int|u_N-u_N^{(R/4)}|\,d\mu_{u_j},
 \qquad u^{(L)}=\max(u,-L).
\end{equation}
By \MathEqref{eq:uniform-energy-test}, the first integral is at most
$C_{q,M}I(u_j,u_N)^\alpha$.
For fixed $N$, the second is at most
$C_{q,M}I(u_N,u_N^{(R/4)})^\alpha$, uniformly
in $j$. The latter discrepancy tends to zero by \MathEqref{eq:clipping-strong}.
Choose $N$ large using \MathEqref{eq:strong-cauchy}, then let $R$ grow, and finally
make the Cauchy error tend to zero. The finitely many indices
$j<N$ have vanishing tails because their first energies are
finite. This proves the first line of \MathEqref{eq:uniform-deep-tails}. The tail form of
\MathEqref{eq:capacity-tail-energy},
\begin{equation}\label{eq:strong-capacity-tail}
 \int_R^\infty t^q C_S(\{u_j<-t\};V)\,dt
       \leq 2^{q+1}\int(-u_j-R/2)^+\,d\mu_{u_j},
\end{equation}
proves the second. For a finite sum of rows, the sublevel union
bound gives the same uniform tail conclusion, with the factor
$(\#\text{slots})^{q+1}$ and a fixed rescaling of $R$.

For \MathEqref{eq:cauchy-weighted-measures}, all slot energies are uniformly bounded. Write
$M_j=M(u_{1,j},\ldots,u_{q,j})$. Positivity of the expansion
of $\mu_{\sum_{i=1}^q u_{i,j}}$ and \MathEqref{eq:uniform-energy-test} give
\begin{equation}\label{eq:weight-freeze}
             \int|u_{0,j}-u_{0,N}|\,dM_j
                  \leq CI(u_{0,j},u_{0,N})^\alpha,
\end{equation}
where $C$ depends on the uniform energies and $q$, but not on
the derivative-row indices. Fix a sufficiently late $N$.
For $\epsilon>0$, the set
$K_\epsilon=\overline{\{u_{0,N}<-\epsilon\}}$ is compact.
On its complement $u_{0,N}=u_{0,N}^{(\epsilon)}$, and hence
\begin{equation}\label{eq:boundary-tightness}
 \begin{split}
 \int_{V\setminus K_\epsilon}(-u_{0,j})\,dM_j
 &\leq \int|u_{0,j}-u_{0,N}|\,dM_j
          +\int(-u_{0,N}^{(\epsilon)})\,dM_j\\
 &\leq CI(u_{0,j},u_{0,N})^\alpha
       +e_1(u_{0,N}^{(\epsilon)})^{1/d}
                          \prod_{i=1}^q e_1(u_{i,j})^{1/d}.
 \end{split}
\end{equation}
The last line is the global constant-one H\"older inequality.
The small-value energy vanishes:
\begin{equation}\label{eq:small-value-energy}
 e_1(u_{0,N}^{(\epsilon)})
   \leq C_q\int_0^\epsilon t^q
                    C_S(\{u_{0,N}<-t\};V)\,dt\longrightarrow0,
\end{equation}
by the power-energy/capacity equivalence. First make the
Cauchy error small, then choose $\epsilon$ small. The
finitely many omitted weights $u_{0,j}$, with derivative
rows still varying, are handled by the same argument with
zero first term: outside their own compact negative sublevel,
the H\"older bound with $u_{0,j}^{(\epsilon)}$ tends uniformly
to zero. A finite union of these compacts proves uniform
tightness. This observation also proves the assertion for
independent row indices, without taking a diagonal assumption
on any derivative slot.
\end{proof}

\subsection{The Borel class and its intrinsic original currents}

\begin{definition}[Strong relative first-energy completion]
\label[definition]{def:r11-strong-first-energy-completion}
Let $\widehat{\mathcal E}^1_{\rm str}(V)$ consist of actual
Borel functions $u:V\to[-\infty,0]$ which are limits of an
actual strongly Cauchy row in every local truncated outer
$S$-capacity. Equality means local outer-capacity equality.
The capacities may equivalently be taken on all compact
cores in the fixed domain $V$, by the analytic capacity-domain comparison in \cref{thm:analytic-estimates}. A choice of actual
defining row is not part of the element of this quotient.
\end{definition}

\begin{proposition}[Borel limits of all strongly Cauchy actual rows]
\label[proposition]{prop:r11-strong-cauchy-borel-limit}
Every row satisfying \MathEqref{eq:strong-cauchy} has a limit in
\cref{def:r11-strong-first-energy-completion}, unique
modulo local outer-capacity zero. Every member of that space is
finite and continuous locally outside ordinary open sets of
arbitrarily small outer capacity. Its pole set is locally
outer-capacity zero.
\end{proposition}

\begin{proof}
Uniform energy and \MathEqref{eq:energy-capacity} make each fixed bounded truncation
of the row Cauchy in every local outer capacity. Apply the
countable compact-core and integer-level diagonal construction
from the bounded completion. A subsequence has compatible
pointwise limits of all bounded truncations off a locally
outer-null set. Its actual pointwise limsup defines a Borel
$u\leq0$ on all of $V$, and the Cauchy property gives the
required capacity convergence for the full row.

There is a uniform deep sublevel estimate. The no-loss
sublevel inequality and the first-energy bound imply
\begin{equation}\label{eq:deep-sublevels}
            C_S(\{u_j<-R\};V)\leq C_q M R^{-d}.
\end{equation}
These sets are ordinary open, so the same bound holds for
their outer capacities. On a fixed compact core, the set
$\{u<-2R\}$ is contained in $\{u_j<-R\}$ together with
a positive gap of the fixed truncations at $3R$.
Pass $j\to\infty$ in the outer-capacity subadditivity bound.
It follows that the pole of $u$ has outer capacity zero on
that core. Outside an ordinary open cover of a sufficiently
deep sublevel, $u$ agrees with one bounded truncation.
Its bounded-completion quasicontinuity then proves the stated
finite local quasicontinuity. The level-set triangle inequality
proves uniqueness. 
\end{proof}

\begin{theorem}[Canonical critical calculus in the strong completion]
\label[theorem]{thm:r11-strong-completed-calculus}
For a finite tuple $\mathbf u$ in
$\widehat{\mathcal E}^1_{\rm str}(V)$, with actual strongly
Cauchy defining rows $\mathbf u_j$, there are intrinsic
coefficient-Radon original currents
\begin{equation}\label{eq:completed-currents}
 \widehat N_\gamma(\mathbf u),\qquad
 \widehat{\mathcal J}'_{\gamma,i},\quad
 \widehat{\mathcal J}''_{\gamma,i},\qquad
 \widehat G^\gamma_{ih}.
\end{equation}
Curvatures have $k=|\gamma|\leq q$; first-order and energy
currents have $k\leq q-1$. All energy degrees are included,
since $k+2\leq q+1=d$. Curvatures are separately closed,
symmetric and middle-positive; energy matrices are positive.
First-order and energy currents are not asserted to be closed.

For any continuous $F$ on the nonpositive orthant satisfying
$|F(x)|\leq A(1-\sum_i x_i)^a$, the actual multiplier limit
\begin{equation}\label{eq:critical-completed-weights}
                  F(\mathbf u_j)P_j
                         \longrightarrow F(\mathbf u)\widehat P
       \qquad(a+b\leq d)
\end{equation}
holds on every original compact smooth complementary test and
every actual moment refinement. Here the degree budget $b$ is
$k$, $k+1$, or $k+2$ for curvature, first-order, or energy.
All these multipliers are coefficient-integrable. The currents
and weighted limits are independent of defining rows and of
outer-capacity-null representatives, and they annihilate the
joint pole set. They satisfy all original recursive identities,
including
\begin{equation}\label{eq:completed-recursion}
 \begin{split}
 D(u_i\widehat N_\gamma)&=\widehat N_{\gamma+e_i},\\
 d'(u_i\widehat N_\gamma)&=\widehat{\mathcal J}'_{\gamma,i},\qquad
 d''(u_i\widehat N_\gamma)=\widehat{\mathcal J}''_{\gamma,i}
                         \qquad(k\leq q-1),
 \end{split}
\end{equation}
and the actual three-term energy identity obtained by placing hats on \MathEqref{eq:three-term-energy}.

For $q+1$ completed slots, the global mixed energy is finite and
\begin{equation}\label{eq:global-completed-energy}
 B(u_{0,j},\ldots,u_{q,j})
      \longrightarrow
 \widehat B(u_0,\ldots,u_q)
       :=\int_V(-u_0)\widehat N(u_1,\ldots,u_q).
\end{equation}
The corresponding positive weighted measures converge against
every bounded original smooth scalar test, including those
without compact support. The assertion also permits independent
indices tending to infinity in the defining rows.
\end{theorem}

\begin{proof}
For a finite tuple define
\begin{equation}\label{eq:joint-tail-function}
 \Phi(R)=\sup_j\int_R^\infty t^q
          C_S\!\left(\left\{\sum_i u_{i,j}<-t\right\};V\right)dt.
\end{equation}
By \cref{lem:r11-strong-cauchy-uniform-tails},
$\Phi(0)<\infty$ and $\Phi(R)\to0$; the same is true for
the supremum over independently large row indices, with the
finitely many earlier indices included if needed. At every
original compact test stage, the actual truncation estimate \MathEqref{eq:graded-tail} give
\begin{equation}\label{eq:truncation-error}
 \|\mu_{P_j-P_j^{(L)},\eta}\|_{\rm TV}
             \leq C_\eta L^{b-d}\Phi(L/4).
\end{equation}
The bounded analytic completion supplies intrinsic currents
$P^{(L)}$ of the completed clipped tuple. Passing $j$ to
infinity in the difference of two bounded stages shows that
$P^{(L)}$ are Cauchy in local coefficient total variation,
by comparing two truncation levels with the same tail bound. Their limit defines $\widehat P$.
Coefficient/refinement identities, positivity and separate
closedness pass on all original tests.

The limiting coefficients have the common modulus
\begin{equation}\label{eq:coefficient-modulus}
    \inf_{L\geq1}\{L^b s+L^{b-d}\Phi(L/4)\}
                              \longrightarrow0\quad(s\downarrow0).
\end{equation}
Its weak-limit proof keeps $L$ fixed until after ordinary-open
variation lower semicontinuity and Radon outer regularity.
Thus it controls both the defining and limiting coefficients,
including the critical case $b=d$. It annihilates every
locally outer-null set and, in particular, the joint pole set.
The bounded quasicontinuous multiplier argument therefore
passes every bounded continuous function of a fixed clipped
tuple.

For endpoint integrability apply that argument to
$\theta_R(z)\min((1+z)^a,H)$, where
$z=-\sum_i u_i$ and $\theta_R$ is zero below $R$ and one
above $2R$. For fixed $R,H$ this is a continuous function
of one sufficiently large clipped tuple; once a slot exceeds
that clipping level its value is the constant $H$.
When $a=0$, use $\theta_R$ itself. The actual row estimates
bound the integral by
$C R^{a+b-d}\Phi(R/8)$. Passing to the bounded multiplier
limit and then $H\uparrow\infty$ proves the same tail bound
for the limiting positive traces. Variation lower
semicontinuity handles signed first-order coefficients, and
positive diagonal controls handle cross energies.
Adding the bounded-value part gives actual coefficient
integrability. Clipping a general $F$ and then letting its
clipping level tend to infinity proves \MathEqref{eq:critical-completed-weights}. Consequently
all multipliers in \MathEqref{eq:completed-recursion} and \MathEqref{eq:three-term-energy} are defined before
distributional differentiation, and the original row
identities pass to the limit.

For two different defining tuples with the same Borel
capacity limits, their bounded completed stages are identical.
The union of their finitely many row families has the same
vanishing-tail property. The preceding construction therefore
gives identical unbounded currents and actual multipliers.
This proves approximation independence.

Finally, \MathEqref{eq:critical-completed-weights} gives local convergence of \MathEqref{eq:cauchy-weighted-measures}.
The uniform global tightness in
\cref{lem:r11-strong-cauchy-uniform-tails}, together
with the uniform finite mixed-energy bound \MathEqref{eq:holder}, identifies
a finite global limiting measure and gives its total-mass
convergence. This is precisely \MathEqref{eq:global-completed-energy}. Tightness also passes
every bounded original scalar test. The independent-index
assertions follow from the same bounds, or by applying the
argument to every sequence of independently diverging indices.
\end{proof}

\subsection{Intrinsic discrepancy and completeness}

Define on this Borel space
\begin{equation}\label{eq:completed-discrepancy}
 \begin{split}
 \widehat e(u)&=\widehat B(u^{[q+1]}),\\
 \widehat I(u,v)&=\widehat e(u)+\widehat e(v)
       -\widehat B(u,v^{[q]})-\widehat B(v,u^{[q]}).
 \end{split}
\end{equation}

\begin{lemma}[Uniform discrepancy continuity on energy balls]
\label[lemma]{lem:r11-discrepancy-uniform-perturbation}
The completed space is closed under finite sums and nonnegative
scalar multiples. Its mixed energy is symmetric, positive and
multilinear on that cone, and satisfies the constant-one H\"older
inequality. The formulas \MathEqref{eq:uniform-energy-test}, \MathEqref{eq:energy-capacity}, \MathEqref{eq:energy-root} and the sharp
midpoint inequality \MathEqref{eq:midpoint} hold in the completion, with their
actual Borel multipliers and local capacities. In particular,
$\widehat I\geq0$, and it vanishes exactly on equal relative
classes.

If $\widehat e(a),\widehat e(a'),\widehat e(b)\leq M$,
then
\begin{equation}\label{eq:uniform-discrepancy}
 \big|\widehat I(a,b)-\widehat I(a',b)\big|
       \leq L_q M^{1-\alpha}\widehat I(a,a')^\alpha,
 \qquad
 L_q=2(q+1)2^{(q-1)\alpha}3^{d(1-\alpha)}.
\end{equation}
The assertion at $M=0$ means that both sides are zero.
\end{lemma}

\begin{proof}
For two strongly Cauchy defining rows, the sum rows have the
required Borel capacity limit. Expand the discrepancy of two
independently late sum rows into finitely many global mixed
energies. Equation \MathEqref{eq:global-completed-energy} makes each term converge to the
corresponding completed mixed energy; the total limiting
discrepancy is zero. Thus the sum row is strongly Cauchy.
Positive scaling is immediate, and zero scaling means the
zero function even at poles. Passing the actual finite mixed
identities and \MathEqref{eq:holder} through \MathEqref{eq:global-completed-energy} now proves all asserted
cone identities.

The actual inequalities \MathEqref{eq:uniform-energy-test}, \MathEqref{eq:energy-root}, and \MathEqref{eq:midpoint} pass by
global mixed-energy convergence. For \MathEqref{eq:uniform-energy-test}, its left side is
the actual weight $|u-v|$, which has degree-one growth and is
bounded by $(-u)+(-v)$; \MathEqref{eq:critical-completed-weights} and the tightness of these two
critical positive mixed measures pass its global integral.
For \MathEqref{eq:energy-capacity}, use defining rows, their convergent energies,
and the level-set triangle inequality for clipped gaps;
then remove deep sublevels using \MathEqref{eq:deep-sublevels} and decrease the
intermediate gap threshold. This also proves that
$\widehat I=0$ gives equality in all local outer capacities.
The converse follows from intrinsic current and multiplier
independence.

It remains to prove \MathEqref{eq:uniform-discrepancy}. It suffices first to consider
actual potentials. Write $h=a-a'$. Telescoping the three
terms of $I(a,b)$ that depend on $a$, and using full mixed
symmetry, expresses their difference as $2(q+1)$ signed
integrals of $-h$ against positive mixed top products whose
slots lie among $a,a',b$. Every such product is bounded by
$\mu_\psi$, where $\psi=a+a'+b$, by the positive expansion.
Thus
\begin{equation}\label{eq:discrepancy-telescope}
 |I(a,b)-I(a',b)|
                  \leq2(q+1)\int|a-a'|\,d\mu_\psi.
\end{equation}
Minkowski gives $e_1(\psi)\leq3^dM$ and
$e_1(a+a')\leq2^dM$. Apply \MathEqref{eq:uniform-energy-test} to \MathEqref{eq:discrepancy-telescope} to obtain
\MathEqref{eq:uniform-discrepancy}. For completed potentials pass through actual
defining rows whose energies tend to the completed energies;
first allow $M+\epsilon$ and then let $\epsilon\downarrow0$.
The zero-energy case follows separately from H\"older.
\end{proof}

\begin{theorem}[Completeness and the intrinsic strong topology]
\label[theorem]{thm:r11-complete-strong-first-energy-space}
Every $\widehat I$-Cauchy sequence in
$\widehat{\mathcal E}^1_{\rm str}(V)$ converges to a unique
relative class in that space, with discrepancy tending to zero.
The actual classical class is strongly dense.

For members $u_j,u$ of this space, one has
\begin{equation}\label{eq:completed-topology}
 \widehat I(u_j,u)\longrightarrow0
 \quad\Longleftrightarrow\quad
 \left\{
 \begin{array}{l}
 u_j\longrightarrow u\text{ in all local truncated outer capacities},\\
 \widehat e(u_j)\longrightarrow\widehat e(u).
 \end{array}\right.
\end{equation}
Strongly convergent finite tuples have every original local
weighted-current limit in \MathEqref{eq:critical-completed-weights}, including all its critical
endpoints, and the global mixed-energy and weighted-measure
limits \MathEqref{eq:global-completed-energy}. The weighted limits here are original-current limits.

\end{theorem}

\begin{proof}
If $(v_l)$ is an actual defining row for a completed $u$, then
for fixed $l$, \MathEqref{eq:global-completed-energy} gives
$\widehat I(v_l,u)=\lim_{r\to\infty}I(v_l,v_r)$.
Its Cauchy property therefore implies
$\widehat I(v_l,u)\to0$. This proves strong density.

Now let $(u_j)$ be completed and $\widehat I$-Cauchy.
The energy-root estimate bounds its energies. Choose actual
$v_j$ with $\widehat I(v_j,u_j)<2^{-j}$ and uniformly bounded
energies. Applying \MathEqref{eq:uniform-discrepancy} in its two slots shows
\begin{equation}\label{eq:diagonal-cauchy}
 I(v_j,v_l)\leq\widehat I(u_j,u_l)
                   +C\bigl(2^{-j\alpha}+2^{-l\alpha}\bigr).
\end{equation}
Hence $(v_j)$ is strongly Cauchy and has a completed Borel
limit $u$. Its defining property and strong density imply
$\widehat I(v_j,u)\to0$. A further use of \MathEqref{eq:uniform-discrepancy} gives
$\widehat I(u_j,u)\to0$. Separation proves uniqueness.
This proves completeness using the uniform discrepancy modulus.

The forward implication of \MathEqref{eq:completed-topology} follows from the completed
energy-root and capacity estimates. For the reverse, the
energies are uniformly bounded. Bounded completed stages and
the strict subcritical truncation bound for top unweighted
currents give local convergence of their top measures. More
explicitly, the stage error has order at most $C M/L$ for
degree $q<d$; first fix $L$, use bounded-completion capacity
continuity, and then let $L$ grow. The same uniform coefficient
modulus permits bounded quasicontinuous weight passage.
Consequently, by first truncating the weight and inserting a
compact nonnegative cutoff, then increasing both,
\begin{equation}\label{eq:cross-energy-liminf}
 \liminf_j\widehat B(u_j,u^{[q]})\geq\widehat e(u),\qquad
 \liminf_j\widehat B(u,u_j^{[q]})\geq\widehat e(u).
\end{equation}
H\"older gives upper bounds tending to $\widehat e(u)$ for
both cross energies. Expanding \MathEqref{eq:completed-discrepancy} proves the converse.
This argument uses no critical weighted continuity to prove
that very continuity.

For a strongly convergent completed tuple choose, at the $j$th
row, classical approximants with discrepancy at most $2^{-j}$,
and sufficiently accurate values on the first $j$ compact
test stages and bounded clipped weights. The defining-row
convergence allows these finitely many requirements
simultaneously. By \MathEqref{eq:uniform-discrepancy}, these actual rows are strongly
Cauchy, with the same Borel limits and energies. The preceding
construction and \MathEqref{eq:critical-completed-weights}--\MathEqref{eq:global-completed-energy} give the asserted limits.
Equivalently, any failure on a fixed test would furnish a
subsequence and actual approximants contradicting those
already proved actual-row limits. This last argument also
covers bounded original scalar tests in the global statement.

\end{proof}


\subsection{Monotone closure and the capacity integral}

\begin{proposition}[Monotone closure]\label[proposition]{prop:completed-monotone}
For completed potentials $a\le b$ relatively quasi-everywhere,
\[
             0\le\widehat I(a,b)\le\widehat e(a)-\widehat e(b).
\]
An increasing sequence in the strong cone converges strongly to its
pointwise monotone Borel limit. The same holds for a decreasing
sequence with uniformly bounded energies. Equalities and monotonicity
are understood outside one locally outer-null set.
\end{proposition}

\begin{proof}
The mixed energy is decreasing in each potential for pointwise order:
this follows from positivity in the weight slot and symmetry in every
other slot. Both cross terms in the discrepancy are consequently at
least $\widehat e(b)$, proving the inequality. Along an increasing
sequence the energies decrease and are nonnegative; along a decreasing
sequence they increase and are bounded by hypothesis. In either case
they converge, so the displayed inequality makes the row strongly
Cauchy. Completeness supplies a limit. A summable-capacity diagonal
subsequence of its bounded truncations converges pointwise off a
locally outer-null set; monotonicity identifies this limit with the
specified Borel monotone limit. Thus the full sequence converges
strongly.
\end{proof}

We also use the following precise capacity characterization from
\cite[Theorem 37.2]{Energy}. A Borel function $u\le0$ represents a
member of the strong cone if and only if each $u^{(m)}$ is a bounded
capacitary limit of actual classical functions on the same domain $V$
with values in $[-m,0]$, and
\begin{equation}\label{eq:completed-capacity-integral}
 J_1^*(u):=\int_0^\infty t^qC_S^*(\{u<-t\};V)\,dt<\infty.
 \qquad
 2^{-d}J_1^*(u)\le\widehat e(u)\le2^dJ_1^*(u).
\end{equation}
The admissibility of the bounded stages is part of this criterion.
The completed sublevel theorem in \cite[Section 36]{Energy} proves the
capacity inequalities and the forward direction. For the reverse,
form $v_j=\max(u,j\rho)=\max(u^{(j)},j\rho)$. Bounded classical
approximants of $u^{(j)}$, followed by this maximum, are bounded below
by the actual first-energy barrier $j\rho$. The fixed-barrier
critical convergence of \cref{thm:analytic-estimates} makes their
limit a strong member $v_j$. The sublevel bounds give
$\widehat e(v_j)\le2^dJ_1^*(u)$. The $v_j$ decrease, hence converge
strongly by \cref{prop:completed-monotone}. On a compact core and for
fixed $R$, strict negativity of $\rho$ gives
$v_j^{(R)}=u^{(R)}$ for all large $j$. The completed limit is therefore
the given capacity class of $u$. The sublevel theorem then gives both
inequalities in \MathEqref{eq:completed-capacity-integral}.

% END sections/08_strong_completion.tex

% BEGIN sections/09_variational_source.tex
\section{The energy dual and the source equation}
\label{sec:variational-source}

In this section $\Ccone=\widehat{\mathcal E}^{1}_{\rm str}(V)$,
and $B,e,I$ denote the completed mixed energy, energy and discrepancy.
Write $r(u)=e(u)^{1/d}$. A difference is formed on the common finite
locus of its two Borel potentials and is assigned any finite value on
the outer-null exceptional set.

\subsection{The Banach space of energy differences}

\begin{definition}[Energy-difference space and its norm]
\label[definition]{def:r12-energy-difference}
Let $\mathbb D_S=\mathscr C-\mathscr C$ be the real vector space of
such Borel differences modulo local outer-capacity equality. Define
\begin{equation}\label{eq:dual-ed3}
 \|h\|_{\rm ed}
   =\inf\{r(a)+r(b):a,b\in\mathscr C,\ h=a-b\}.
\end{equation}

\end{definition}

\begin{lemma}[Continuous extension of the mixed energy]
\label[lemma]{lem:r12-difference-mixed-extension}
There is a unique symmetric real $d$-linear form on $\mathbb D_S$
extending $B$, and it satisfies
\begin{equation}\label{eq:dual-ed4}
       |B(h_0,\ldots,h_q)|\leq\prod_{i=0}^q\|h_i\|_{\rm ed}.
\end{equation}
Moreover $\|u\|_{\rm ed}=r(u)$ for every $u\in\mathscr C$.
\end{lemma}

\begin{proof}
Expand each difference into its two cone terms and extend by the
finite multilinear expansion. If $a-b=a'-b'$, then
$a+b'=a'+b$ in the relative cone. Separate additivity of
the completed $B$ therefore makes the expansion independent of that
presentation, one slot at a time. Positive real homogeneity and
changing the signs of differences give real multilinearity.
For specified decompositions $h_i=a_i-b_i$, the triangle inequality
and cone H\"older give the bound $\prod_i(r(a_i)+r(b_i))$.
Taking independent infima proves \MathEqref{eq:dual-ed4}.

The definition gives $\|u\|_{\rm ed}\leq r(u)$. Applying \MathEqref{eq:dual-ed4}
to the diagonal tuple gives $e(u)\leq\|u\|_{\rm ed}^{d}$, so the
reverse inequality follows. This argument also covers zero energy.
\end{proof}

\begin{theorem}[Actual Borel realization and Banach completeness]
\label[theorem]{thm:r12-energy-difference-Banach}
Formula \MathEqref{eq:dual-ed3} is a norm, $\mathbb D_S$ is complete for that norm,
and $\mathscr C$ is a closed convex cone in this Banach space.
Norm convergence implies convergence in the full local outer
$S$-capacity topology. More precisely, with
$C_0=d2^d$, one has
\begin{equation}\label{eq:dual-ed5}
 C_S^*(\{|h|>t\};V)\leq C_0 t^{-d}\|h\|_{\rm ed}^{d}
                    \qquad(h\in\mathbb D_S,\ t>0).
\end{equation}
For $u,v\in\mathscr C$,
\begin{equation}\label{eq:dual-ed6}
 I(u,v)\leq\bigl(r(u)^q+r(v)^q\bigr)\|u-v\|_{\rm ed}.
\end{equation}

\end{theorem}

\begin{proof}
The infimum in \MathEqref{eq:dual-ed3} is finite. Minkowski on $\mathscr C$ proves
the triangle inequality; switching the two cone terms proves
symmetry, and homogeneity follows by rescaling decompositions.
The capacity integral bound \MathEqref{eq:completed-capacity-integral} and
monotonicity of the sublevels give
\[
 \frac{t^d}{d}C_S^*(\{u<-t\};V)
       \leq\int_0^t s^{d-1}C_S^*(\{u<-s\};V)\,ds
       \leq2^d e(u).
\]
For $h=a-b$, the set $\{|h|>t\}$ is contained, modulo an
outer-null set, in $\{a<-t\}\cup\{b<-t\}$. Hence its outer
capacity is at most $C_0t^{-d}(r(a)^d+r(b)^d)$, bounded by
$C_0t^{-d}(r(a)+r(b))^d$. Infimizing proves \MathEqref{eq:dual-ed5}, and shows
that a zero seminorm has the zero relative Borel class. Local
outer-null sets here also have zero outer capacity in the fixed
domain: use a countable compact exhaustion and the analytic domain comparison on its cores. The same comparison transfers norm
convergence to every smaller-domain original capacity.

We prove completeness inside the space of actual Borel differences,
rather than introducing an abstract completion with unidentified
elements. For a norm-Cauchy sequence choose a subsequence $(h_j)$
whose successive increments have decompositions
\begin{equation}\label{eq:dual-ed7}
 h_{j+1}-h_j=a_j-b_j,\qquad
 \sum_{j\geq1}(r(a_j)+r(b_j))<\infty.
\end{equation}
Write $h_1=a_0-b_0$. The partial sums
$A_N=a_0+\sum_{j=1}^N a_j$ and
$B_N=b_0+\sum_{j=1}^N b_j$ decrease as Borel potentials,
outside one countable union of outer-null sets. Their energy
roots are bounded by the corresponding sums of roots, so the
\cref{prop:completed-monotone} gives strong limits
$A,B\in\mathscr C$. Apply the same theorem to each tail series.
It yields $A-A_N,B-B_N\in\mathscr C$ and
\begin{equation}\label{eq:dual-ed8}
 r(A-A_N)+r(B-B_N)
          \leq\sum_{j>N}(r(a_j)+r(b_j))\longrightarrow0.
\end{equation}
The identities with the tails hold on the common finite locus,
hence in the relative Borel quotient. Thus $h=A-B\in\mathbb D_S$
is the norm limit of the subsequence and of the original Cauchy
sequence. No abstract vector invisible to Borel capacity has been
adjoined.

For $u,v\in\mathscr C$, multilinearity gives
\[
 I(u,v)=B(u-v,u^{[q]})-B(u-v,v^{[q]}).
\]
Apply \MathEqref{eq:dual-ed4} to obtain \MathEqref{eq:dual-ed6}. If a sequence in $\mathscr C$
converges in norm to $h\in\mathbb D_S$, its energy roots are
bounded by the norm triangle inequality and the preceding lemma.
Equation \MathEqref{eq:dual-ed6} makes it strongly Cauchy. \cref{thm:r11-complete-strong-first-energy-space} supplies
a limit $u\in\mathscr C$. The norm limit and the strong limit
have the same local capacity class by \MathEqref{eq:dual-ed5} and \MathEqref{eq:completed-topology}. Thus $h=u$, proving norm closedness.
Convexity is the completed cone addition property.
\end{proof}

\begin{definition}[Curvature as an energy-dual functional]
\label[definition]{def:r12-dual-curvature}
For $u\in\mathscr C$ put
\begin{equation}\label{eq:dual-ed9}
       J(u)(h)=B(h,u^{[q]})=-\int_Vh\,d\widehat\mu_u
                      \qquad(h\in\mathbb D_S).
\end{equation}
The displayed integral is literal and absolutely convergent: for
$h=a-b$, $|h|\leq(-a)+(-b)$ off a common null set, and the
two mixed first energies are finite. The current measure ignores
that exceptional set. Thus $J(u)$ is a genuine functional in
$\mathbb D_S^*$, with its original current interpretation.
\end{definition}

\begin{theorem}[Exact dual norm and strong continuity of curvature]
\label[theorem]{thm:r12-dual-curvature-continuity}
For every continuous linear functional $\ell$ on $\mathbb D_S$,
\begin{equation}\label{eq:dual-ed10}
 \|\ell\|_*=
     \sup_{\substack{a\in\mathscr C\\r(a)\leq1}}|\ell(a)|.
\end{equation}
For $u\in\mathscr C$ one has $\|J(u)\|_*=r(u)^q$.
For $u,v\in\mathscr C$ and $R=\max(r(u),r(v))$,
\begin{equation}\label{eq:dual-ed11}
 \begin{split}
 \|J(u)-J(v)\|_*&\leq qR^{q-1}\|u-v\|_{\rm ed},\\
 \|J(u)-J(v)\|_*&\leq\sqrt q\,R^{(q-1)/2}I(u,v)^{1/2}.
 \end{split}
\end{equation}
At $R=0$ both sides mean zero; for $q=1$ the power $R^0$ is
one. In particular strong convergence of potentials gives norm
convergence of their source functionals, in every remaining degree.
\end{theorem}

\begin{proof}
If the supremum in \MathEqref{eq:dual-ed10} is $L$, a decomposition $h=a-b$
gives $|\ell(h)|\leq L(r(a)+r(b))$. Taking the infimum yields
$\|\ell\|_*\leq L$; the reverse follows from
$\|a\|_{\rm ed}=r(a)$. Equation \MathEqref{eq:dual-ed4} bounds the norm of
$J(u)$ by $r(u)^q$. For $r(u)>0$, evaluation on $u/r(u)$
gives equality; if $r(u)=0$, the class and functional are zero.

The finite telescoping expansion gives
\begin{equation}\label{eq:dual-ed12}
 (J(u)-J(v))(h)
    =\sum_{l=0}^{q-1}B(h,u-v,u^{[l]},v^{[q-1-l]}).
\end{equation}
Equation \MathEqref{eq:dual-ed4} proves the first bound in \MathEqref{eq:dual-ed11}.
For the sharper strong-topology bound define
\[
 H(f,g)=\sum_{l=0}^{q-1}
                 B(f,g,u^{[l]},v^{[q-1-l]}).
\]
This is a positive semidefinite bilinear form on $\mathbb D_S$.
Indeed each summand is the integrated energy form with the
fixed completed positive closed mixed background. Positivity on
finite differences follows from \MathEqref{eq:global-energy-matrix}, passed through the global completed-energy convergence \MathEqref{eq:global-completed-energy}.  We have
\[
 H(u-v,u-v)=I(u,v),\qquad
 H(a,a)\leq q r(a)^2R^{q-1}\quad(a\in\mathscr C).
\]
Cauchy--Schwarz for this positive semidefinite form and \MathEqref{eq:dual-ed12}
bound the functional on every $a$ with $r(a)\leq1$.
Equation \MathEqref{eq:dual-ed10} proves the second bound. Strongly convergent
rows have bounded energy roots by \MathEqref{eq:energy-root}, so the final assertion
follows. 
\end{proof}


\subsection{Closed sources and variational normals}

\begin{theorem}[A complete curvature-dual metric and closed source image]
\label[theorem]{thm:r12-closed-full-source-image}
The image $J(\mathscr C)$ is norm closed in $\mathbb D_S^*$.
The formula
\begin{equation}\label{eq:dual-ni1}
                 d_J(u,v)=\|J(u)-J(v)\|_*
\end{equation}
defines a complete metric on $\mathscr C$, inducing precisely the strong discrepancy topology. For $u,v\in\mathscr C$,
\begin{equation}\label{eq:dual-ni2}
 \begin{split}
 I(u,v)&\leq(r(u)+r(v))d_J(u,v),\\
 d_J(u,v)&\leq\sqrt q\,
     \max(r(u),r(v))^{(q-1)/2}I(u,v)^{1/2}.
 \end{split}
\end{equation}

\end{theorem}

\begin{proof}
Multilinearity gives the exact identity
$I(u,v)=(J(u)-J(v))(u-v)$. The dual inequality and
$\|u-v\|_{\rm ed}\leq r(u)+r(v)$ prove the first estimate.
The second is \MathEqref{eq:dual-ed11}. Thus $d_J=0$ gives $I=0$ and equal
relative classes. Symmetry and the triangle inequality come
from the dual norm, not from an asserted triangle inequality
for $I$.

For completeness, mixed H\"older also gives the explicit root bound
\[
 I(u,v)\geq\bigl(r(u)-r(v)\bigr)
                   \bigl(r(u)^q-r(v)^q\bigr)
                \geq|r(u)-r(v)|^{q+1}.
\]
The first inequality follows by bounding the two mixed terms in
$e(u)+e(v)-B(u,v^{[q]})-B(v,u^{[q]})$ separately. The second
is the elementary inequality $x^q-y^q\geq(x-y)^q$ for
$x\geq y\geq0$. Hence convergence to a fixed potential in
$I$ automatically bounds the energy roots.

Suppose $J(u_j)$ is norm Cauchy. Its norms are bounded, so
$r(u_j)^q=\|J(u_j)\|_*$ bounds the energy roots. The first
estimate in \MathEqref{eq:dual-ni2} makes $(u_j)$ strongly Cauchy. \cref{thm:r11-complete-strong-first-energy-space} gives $u\in\mathscr C$ with $I(u_j,u)\to0$.
Its energy roots remain bounded, and the second estimate gives
$J(u_j)\to J(u)$ in norm. Every norm limit of image elements
is consequently in the image, proving closedness and metric
completeness. The same inequalities prove equivalence of the
two convergences: strong convergence bounds the roots by the
\MathEqref{eq:energy-root}, whereas $d_J$ convergence bounds
them by the exact dual norm identity. These arguments also
identify their Cauchy sequences.
\end{proof}

\begin{definition}[The strong-continuous energy dual]
\label[definition]{def:r12-strong-continuous-dual}
Let $\mathbb D_{S,\mathrm{s}}^*$ consist of those
$\ell\in\mathbb D_S^*$ for which
\begin{equation}\label{eq:dual-ni4}
  u_j,u\in\mathscr C,\ I(u_j,u)\to0
                 \quad\Longrightarrow\quad\ell(u_j)\to\ell(u).
\end{equation}

\end{definition}

\begin{lemma}[A closed dual class containing every curvature]
\label[lemma]{lem:r12-strong-dual-closed}
The class $\mathbb D_{S,\mathrm{s}}^*$ is a norm-closed linear
subspace of $\mathbb D_S^*$, and $J(\mathscr C)$ is contained
in it. Consequently the norm-closed linear span of all curvature
functionals is contained in this class.
\end{lemma}

\begin{proof}
Linearity is immediate. A strongly convergent row has bounded
$r(u_j)$, hence is bounded in $\mathbb D_S$. If $\ell_m\to\ell$
in norm, the error $|(\ell-\ell_m)(u_j-u)|$ is uniformly small
for large $m$. Fixing $m$ then invokes \MathEqref{eq:dual-ni4}, proving closure.
For fixed $v\in\mathscr C$, the completed tester estimate \MathEqref{eq:uniform-energy-test} bounds $|J(v)(u_j-u)|$ by
$\int|u_j-u|\,d\widehat\mu_v\to0$. Its constants are finite
on the bounded-energy family. This proves the curvature assertion.
\end{proof}

Let
\begin{equation}\label{eq:dual-ni5}
 \mathscr C^\circ
   =\{n\in\mathbb D_S^*:n(v)\leq0\text{ for every }v\in\mathscr C\}.
\end{equation}


\begin{theorem}[All-degree variational normal decomposition]
\label[theorem]{thm:r12-nonlinear-normal-decomposition}
For every $\ell\in\mathbb D_{S,\mathrm{s}}^*$, the functional
\begin{equation}\label{eq:dual-ni6}
                F_\ell(v)=\ell(v)-e(v)/d\qquad(v\in\mathscr C)
\end{equation}
has a unique maximizer $u=\Pi(\ell)$. Every maximizing sequence
converges strongly to $u$. Equivalently, there is a unique pair
$(u,n)$ with
\begin{equation}\label{eq:dual-ni7}
       \ell=J(u)+n,\qquad
       u\in\mathscr C,\quad n\in\mathscr C^\circ,
                           \quad n(u)=0.
\end{equation}
The residual also lies in $\mathbb D_{S,\mathrm{s}}^*$.
Define the nonnegative support gauge
\begin{equation}\label{eq:dual-ni8}
       p(\ell)=\sup_{\substack{v\in\mathscr C\\r(v)\leq1}}\ell(v).
\end{equation}
Then
\begin{equation}\label{eq:dual-ni9}
 \begin{split}
 e(u)&=\ell(u),\qquad r(u)^q=p(\ell)\leq\|\ell\|_*,\\
 \sup F_\ell&=\frac qd\,p(\ell)^{d/q}.
 \end{split}
\end{equation}
If $\ell$ is nonnegative on $\mathscr C$, then
$p(\ell)=\|\ell\|_*$. 
\end{theorem}

\begin{proof}
Since $|\ell(v)|\leq\|\ell\|_*r(v)$,
$F_\ell(v)\leq\|\ell\|_*r(v)-r(v)^d/d$. Thus its supremum
$L\geq F_\ell(0)=0$ is finite, and every maximizing sequence
has bounded energies. In the midpoint gap the linear term
cancels. \MathEqref{eq:midpoint} gives
\[
 c_q I(v_j,v_k)
     \leq L-\tfrac12(F_\ell(v_j)+F_\ell(v_k))\longrightarrow0,
 \quad c_q=\frac{1-2^{-q}}{2d}.
\]
Strong completeness provides $u\in\mathscr C$. Its energy
converges, and \MathEqref{eq:dual-ni4} passes the datum term, so $F_\ell(u)=L$.
The same midpoint estimate proves uniqueness and convergence
of every maximizing sequence. Actual classical strong density
also gives the same supremum over the actual classical class.

The continuous $d$-linear extension $B$ makes $e/d$ a polynomial
on $\mathbb D_S$ whose differential at $u$ is $J(u)$.
Only segments inside $\mathscr C$ are used for maximization.
Their first variation gives
\begin{equation}\label{eq:dual-ni10}
                   (\ell-J(u))(v-u)\leq0\qquad(v\in\mathscr C).
\end{equation}
Conversely convexity of $e$ on the cone makes this sufficient.
Taking the competitors $0,2u$ gives $(\ell-J(u))(u)=0$.
Then \MathEqref{eq:dual-ni10} is precisely the polar condition in \MathEqref{eq:dual-ni7}.
Conversely \MathEqref{eq:dual-ni7} implies \MathEqref{eq:dual-ni10}. The pair is unique because
the maximizer is unique and $n=\ell-J(u)$. Its strong continuity
follows from the preceding lemma.

Scaling along $tu$ gives $e(u)=\ell(u)$, also when $u=0$.
For each $v\in\mathscr C$, the polar inequality gives
$\ell(v)\leq J(u)(v)\leq r(v)r(u)^q$. Thus
$p(\ell)\leq r(u)^q$. If $u\ne0$, evaluation on $u/r(u)$
gives the reverse inequality; if $u=0$, the polar condition
and the allowed zero tester give $p(\ell)=0$. This proves
\MathEqref{eq:dual-ni9}. Finally \MathEqref{eq:dual-ed10} says
$\|\ell\|_* =\max(p(\ell),p(-\ell))$; if $\ell\geq0$ on
the cone, the latter second term is zero.
\end{proof}

\begin{theorem}[A quantitative nonlinear retraction of source functionals]
\label[theorem]{thm:r12-source-retraction}
Put $\mathcal P(\ell)=J(\Pi(\ell))$ on
$\mathbb D_{S,\mathrm{s}}^*$. This is a retraction onto the norm-closed
image $J(\mathscr C)$, and is positively homogeneous of degree one.
For $\|\ell\|_*,\|m\|_*\leq R$, with $R>0$, one has
\begin{equation}\label{eq:dual-ni11}
 \begin{split}
 I(\Pi(\ell),\Pi(m))&\leq2R^{1/q}\|\ell-m\|_*,\\
 \|\mathcal P(\ell)-\mathcal P(m)\|_*
                         &\leq\sqrt{2qR}\,\|\ell-m\|_*^{1/2}.
 \end{split}
\end{equation}
The first inequality has the more precise right side
$(\ell-m)(\Pi(\ell)-\Pi(m))$. If
$n=\ell-\mathcal P(\ell)$ and $R=\|\ell\|_*>0$, then
\begin{equation}\label{eq:dual-ni12}
 \left(\sqrt{qR+\|n\|_*}-\sqrt{qR}\right)^2
 \leq\operatorname{dist}_*(\ell,J(\mathscr C))\leq\|n\|_*.
\end{equation}
In particular a nonzero residual is a quantitative obstruction to
source-image approximation, not merely a failure of one chosen
maximization algorithm.
\end{theorem}

\begin{proof}
If $\ell=J(a)$, the pair $(a,0)$ satisfies \MathEqref{eq:dual-ni7}, so uniqueness
gives the retraction. Homogeneity of $J$ gives
$\Pi(t\ell)=t^{1/q}\Pi(\ell)$ and
$\mathcal P(t\ell)=t\mathcal P(\ell)$ for $t\geq0$;
the case $t=0$ is the zero class.

Write $u=\Pi(\ell)$ and $v=\Pi(m)$. Use $v,u$ as the
competitors in their two inequalities \MathEqref{eq:dual-ni10}, and add. It yields
$I(u,v)\leq(\ell-m)(u-v)$. Duality bounds this by
$\|\ell-m\|_*(r(u)+r(v))$; \MathEqref{eq:dual-ni9} bounds the roots by
$R^{1/q}$. The second line of \MathEqref{eq:dual-ni11} now follows from \MathEqref{eq:dual-ed11}.
If $R=0$ both projections are zero.

Let $\delta$ be the distance in \MathEqref{eq:dual-ni12}. Since $J(0)=0$,
$\delta\leq R$. Choose $\eta_j\in J(\mathscr C)$ with
$\delta_j=\|\ell-\eta_j\|_*\to\delta$. Their norms are
at most $R+\delta_j$. Retraction and \MathEqref{eq:dual-ni11} imply
\[
 \|n\|_*\leq\delta_j+
               \sqrt{2q(R+\delta_j)\delta_j}.
\]
Passing to the limit and using $\delta\leq R$ gives
$\|n\|_*\leq\delta+2\sqrt{qR\delta}$. Solving this
quadratic inequality in $\sqrt\delta$ gives the lower bound.
The upper bound uses the particular image element $\mathcal P(\ell)$.
\end{proof}


\subsection{One geometric potential and all compact smooth directions}

We next construct a potential that permits both signs of each original
compact smooth scalar direction. The construction is made from the
geometry of $V$ before choosing $S$.

\begin{lemma}[Countable exact moments and compensators]
\label[lemma]{lem:geometric-compensators}
There are compact moment cores $K_l\Subset V$ and smooth classical
psh functions $a_l$ on $V$, with $0\le a_l\le1$, such that every
original compact smooth scalar test can be localized to finitely many
of these cores. On each selected core, $Da_l$ dominates a positive
multiple of the identity form in all the moment directions of the
localized test and its smooth partition. The same assertion holds
after every actual finite refinement.
\end{lemma}

\begin{proof}
We use the exact unit replacement of \cite[Lemma 2.1]{Energy}.
For clarity, its countability comes from a countable dense subring
$A_0$ of the separable affinoid algebra $\mathcal O(B)$. Laurent
inequalities in $A_0$ with rational positive thresholds give a
countable neighborhood system. On a suitable such Laurent domain,
a unit $u$ is approximated by a quotient $b/c$ from $A_0$ so closely
that $\|u-b/c\|\|u^{-1}\|<1$. Thus $|u|=|b/c|$ at every point
of the domain. The logarithmic moment and all its CLD differentials
are exactly preserved. Close this countable collection under finite
refinements and select countable cofinal compact cores.

For a finite moment list $t_i=\log|b_i/c_i|$ on a core, choose a
smooth nonnegative function $\kappa$ of integral one supported in
$(-1,0)$ and set
\[
 \zeta(s)=\int_{-\infty}^{s}(s-r)\kappa(r)\,dr,
 \qquad F(s)=\zeta(s)^2.
\]
The function $F$ is convex and nondecreasing, vanishes for $s\le-1$,
and has $F''=2$ for $s\ge0$. Choose $A$ so that
$\log|b_i|+A$ and $\log|c_i|+A$ are positive near the core.
The sum of $F(\log|b_i|+A)$ and $F(\log|c_i|+A)$ is smooth and
classical on the affinoid interior, with the zero branches extended
by zero near analytic zeros. Its Hessian dominates
$\sum_i d't_i\wedge d''t_i$ near the core, by
$(\xi-\eta)^2\le2\xi^2+2\eta^2$. Normalize this bounded sum on
$\overline V$ to obtain $a_l$. Finite localization of any actual test
uses finitely many exact moment cores and includes the partition
moments. This proves the assertion with the full original test sheaf.
\end{proof}

\begin{proposition}[A simultaneous geometric probe]
\label[proposition]{prop:geometric-probe}
There is a bounded continuous classical potential $\psi_0$ determined
by the box alone such that, for every extending allowed $S$,
\[
 -1<\psi_0<0,\qquad \rho\le\psi_0,\qquad
 \mu_{\psi_0}(V)\le\mu_\rho(V)<\infty.
\]
It has compact negative sublevels and ordinary zero boundary value.
For every original real compact smooth scalar $\phi$, there is
$\epsilon_\phi>0$, independent of $S$, such that
\begin{equation}\label{eq:two-sided-probe}
 \psi_0+t\phi\in\mathcal E^1_{S,\partial}(V),
       \qquad |t|\le\epsilon_\phi.
\end{equation}
These functions are bounded continuous classical psh, lie between
$-1$ and $0$, and have finite total top mass. On a fixed finite
presentation one may choose
\begin{equation}\label{eq:probe-seminorm}
 \epsilon_\phi\ge
       \frac{c_{\mathfrak a}}{\max(1,\|\phi\|_{2,\mathfrak a})},
\end{equation}
where $\|\cdot\|_{2,\mathfrak a}$ is the sum of the real $C^2$ norms
of its representatives on that presentation.
\end{proposition}

\begin{proof}
\textbf{Step 1: construct the potential.}
Choose $A_l\ge1$ with $A_l\rho<-2$ near $K_l$, and set
\[
 \psi_l=\max(a_l-2,A_l\rho),\qquad
 c_l=\frac{2^{-l}}{1+A_l},\qquad
 \psi_0=\sum_{l=1}^\infty c_l\psi_l.
\]
Each $\psi_l$ is bounded between $-2$ and $0$, and
$\psi_l\ge A_l\rho$. The series converges uniformly. Its partial
sums decrease and are classical, so the finite continuous limit is
classical by bounded decreasing closure. Moreover,
\[
 T:=\sum_lc_lA_l=1-\sum_lc_l<1,\qquad
                  \psi_0\ge T\rho> -1.
\]
Strict negativity follows from any summand. The lower barrier gives
compact negative sublevels and zero boundary value. Boundary
comparison with $c\rho$, $c>1$, on the increasing compact-gap sets
$\{c\rho<\psi_0-\epsilon\}$, followed by $\epsilon\downarrow0$
and $c\downarrow1$, gives the mass bound. Thus all finite positive
power energies of $\psi_0$ are finite. Omitting finitely many
summands gives the same properties for the remainder.

\textbf{Step 2: perturb each selected moment core.}
Localize $\phi=\sum_{l\in F}\phi_l$ on finitely many cores from
\cref{lem:geometric-compensators}, combining repeated indices. Near
$\supp\phi_l$, one has $\psi_l=a_l-2$. The Hessian of the actual
smooth representative of $\phi_l$ is bounded in the selected moment
directions. The strict positive Hessian contributed by the corresponding
convex-flat numerator and denominator terms of $a_l$ dominates both
signs of this perturbation. Hence there is $\delta_l>0$ such that
$\psi_l+\tau\phi_l$ is classical for $|\tau|\le\delta_l$.
Outside the support it equals $\psi_l$, so ordinary gluing applies.
Decrease $\delta_l$ to retain strict negativity on its compact
support.

Write $\psi_0=r_F+\sum_{l\in F}c_l\psi_l$ and choose
$\epsilon_\phi/c_l\le\delta_l$ for every selected $l$. Then
\[
 \psi_0+t\phi
       =r_F+\sum_{l\in F}c_l
                      (\psi_l+(t/c_l)\phi_l)
\]
is classical and negative. Decrease $\epsilon_\phi$ again so that
$\epsilon_\phi\|\phi\|_\infty\le(1-T)/2$. This preserves the
lower bound $-1$. The perturbation agrees with $\psi_0$ off one
compact set, so its boundary condition is unchanged. Its top source
agrees there by locality; relative CLN gives finite mass on a larger
compact neighborhood. It therefore has finite total mass and belongs
to every finite positive power class.

For a fixed finite presentation, the Hessian bounds, the negativity
margins and the finitely many coefficients $c_l$ are uniform on the
unit ball of its $C^2$ seminorm. Scaling this estimate gives
\MathEqref{eq:probe-seminorm}.
\end{proof}

In particular every compact original smooth scalar belongs to
$\mathbb D_S$: it is the difference of the two cone members
$\epsilon_\phi^{-1}(\psi_0+\epsilon_\phi\phi)$ and
$\epsilon_\phi^{-1}\psi_0$.

\subsection{The scalar source defect}

A Radon measure $\mu$ defines a \emph{literal energy functional} if
$\ell_\mu(h)=-\int_Vh\,d\mu$ is absolutely defined on $\mathbb D_S$,
independent of relative representatives, and continuous for its norm.
We impose strong continuity when writing
$\ell_\mu\in\mathbb D_{S,\mathrm s}^*$.

\begin{theorem}[One probe detects the literal source equation]
\label[theorem]{thm:single-probe}
Suppose $\ell=J(u)+n$ is a normal pair, with
$u\in\Ccone$, $n(v)\le0$ for every $v\in\Ccone$, and $n(u)=0$.
Set
\[
 \Xi_{\psi_0}(\ell,u)=J(u)(\psi_0)-\ell(\psi_0)=-n(\psi_0)\ge0.
\]
For every original real compact smooth scalar $\phi$,
\begin{equation}\label{eq:single-probe-estimate}
                 |n(\phi)|\le\epsilon_\phi^{-1}
                                      \Xi_{\psi_0}(\ell,u).
\end{equation}
Now let $\mu$ be positive Radon with
$\ell_\mu\in\mathbb D_{S,\mathrm s}^*$, and let
$u=\Pi(\ell_\mu)$. Then the following are equivalent:
\begin{enumerate}
\item $\displaystyle\Xi_{\psi_0}(\mu)
        :=\int_V(-\psi_0)\,d\widehat\mu_u
             -\int_V(-\psi_0)\,d\mu=0$;
\item $\widehat\mu_u=\mu$ on every original compact smooth scalar test;
\item there exists $v\in\Ccone$ with $\widehat\mu_v=\mu$.
\end{enumerate}
The solution in \textup{(iii)} is unique and equals $u$. For all data
in this class,
\begin{equation}\label{eq:literal-defect-bound}
 \left|\int_V\phi\,d(\mu-\widehat\mu_u)\right|
                       \le\epsilon_\phi^{-1}\Xi_{\psi_0}(\mu).
\end{equation}
The same $\psi_0$ works for every extending background, every legal
remaining degree, and every actual refinement of the original tests.
\end{theorem}

\begin{proof}
Apply the normal inequality to the two cone members
$\psi_0\pm\epsilon_\phi\phi$ from
\MathEqref{eq:two-sided-probe}. It gives
$0\ge n(\psi_0)\pm\epsilon_\phi n(\phi)$, which is
\MathEqref{eq:single-probe-estimate}.

For a literal datum both probe integrals are finite by the energy-dual
bound and mixed H\"older. The normal on a smooth test is
$-\int\phi\,d(\mu-\widehat\mu_u)$, proving
\MathEqref{eq:literal-defect-bound}. If $\Xi=0$, all these test
integrals vanish. Original scalar density and Radon uniqueness give
equality of the measures. Conversely, measure equality gives equality
of the two finite probe integrals. Finally, a source potential $v$ has
$J(v)=\ell_\mu$ by literal integration. The normal pair $(v,0)$ and
the variational uniqueness theorem give $v=u$. This proves all three
equivalences.
\end{proof}

For a nonliteral functional the first part detects its compact-test
shadow. The literal hypothesis in the equivalence identifies that
shadow with the specified Radon measure.

\begin{corollary}[Additive strictness and dense admissible sources]
\label[corollary]{cor:additive-strictness}
Under the literal hypotheses of \cref{thm:single-probe}, if
$u=v+\varepsilon\psi_0$ for some $v\in\Ccone$ and
$\varepsilon>0$, then $\widehat\mu_u=\mu$.
For any $a\in\Ccone$, the potentials $a+\varepsilon\psi_0$ admit
both signs of every compact smooth direction, tend to $a$ in the
energy-difference norm, and satisfy
\[
 \|J(a+\varepsilon\psi_0)-J(a)\|_*
       \le(r(a)+\varepsilon r(\psi_0))^q-r(a)^q.
\]
Their sources are therefore dense in $J(\Ccone)$.
\end{corollary}

\begin{proof}
The normal satisfies
$0=n(u)=n(v)+\varepsilon n(\psi_0)$, and both terms are
nonpositive. Hence $\Xi=0$, so \cref{thm:single-probe} applies.
For the second assertion use
$a+\varepsilon(\psi_0+t\phi)\in\Ccone$ and
$\|\varepsilon\psi_0\|_{\rm ed}=\varepsilon r(\psi_0)$.
Expanding the $q$ source slots and bounding each mixed functional by
the product of its energy roots gives the displayed binomial estimate.
\end{proof}

% END sections/09_variational_source.tex

% BEGIN sections/10_linear_dirichlet.tex
\section{The relative linear renormalized Dirichlet problem}
\label{sec:linear-dirichlet}

Here the remaining degree is $q=1$, so $S$ has degree $(n-1,n-1)$.
The ambient dimension is arbitrary. This section uses the following
slightly more general domain: choose one compact original affinoid $B$,
ordinary opens $V\Subset\Omega$ with
$\overline\Omega\subset\operatorname{Int}_X(B)$, and an allowed
background $S$ on $\Omega$. The Laurent domains used above satisfy
these conditions after an intermediate neighborhood is chosen.

\subsection{The original energy form and its diffuse data}

The faithful-form construction
\cite[Theorem \emph{A faithful coercive original form without algebraizability}
and its part-space corollary]{Energy} gives a bounded continuous
classical $0\le h\le1$ on $\Omega$ and a positive reference measure
$\lambda=D(hS)$. It is faithful to the original coefficients:
$\lambda(E)=0$ if and only if $\mathbf1_ES=0$ for Borel $E$.
Put $F=V\cap\supp S$. The measure $\lambda|_F$ is finite.
The energy closure $H$ of the original real smooth compact scalar
core on $V$, modulo zero energy, is a regular strongly local symmetric
Dirichlet form on $F$, with
\begin{equation}\label{eq:linear-form}
 \mathcal E(f,g)=\int_V\tfrac12
        (d'f\wedge d''g+d'g\wedge d''f)\wedge S,
 \qquad \mathcal E(f,f)\ge\tfrac12\|f\|_{L^2(\lambda)}^2.
\end{equation}
For an original compact smooth $\phi$,
\begin{equation}\label{eq:original-form-coefficients}
 D(\phi S)=b_\phi\lambda,\quad b_\phi\in L^\infty_c(\lambda),
 \qquad
 \mathcal E(v,\phi)=-\langle D(vS),\phi\rangle
                   =-\int_Fv b_\phi\,d\lambda\quad(v\in H).
\end{equation}
Every original compact coefficient of $S$ is bounded by a constant
multiple of $\lambda$. Multiplication by a fixed core cutoff obeys
\begin{equation}\label{eq:cutoff-energy}
 \mathcal E(\chi v,\chi v)
 \le2\|\chi\|_\infty^2\mathcal E(v,v)+2C_\chi\|v\|_2^2.
\end{equation}
The faithful measure is obtained from a summable sequence of the exact
moment compensators of \cref{lem:geometric-compensators}. Thus the
form and its current identities retain every original refinement.

Let $c_V$ be the transient zero-capacity of this part form; for
quasi-open sets it is the infimum of $\mathcal E(v,v)$ over $v\in H$
with $\widetilde v\ge1$ quasi-everywhere. A finite signed Radon datum
$\mu$ on $F$ is \emph{energy diffuse} if $|\mu|$ vanishes on every
$c_V$-polar Borel set. The tilde denotes the canonical
energy-quasicontinuous representative. This energy capacity and the
relative Bedford--Taylor capacity play different roles.

For $f\in L^2(F,\lambda)$, let $Gf\in H$ be the Riesz potential,
$\mathcal E(Gf,z)=\int fz\,d\lambda$. The original barrier gives
$|Gf|\le(2-h)\|f\|_\infty\le2\|f\|_\infty$ for bounded $f$.
Set $T_k(s)=\max(-k,\min(k,s))$.

\begin{theorem}[Transposition with the full original current]
\label[theorem]{thm:linear-transposition}
Every finite signed energy-diffuse Radon measure $\mu$ has a unique
$u_\mu\in L^1(F,\lambda)$ satisfying
\begin{equation}\label{eq:transposition}
 \int_Fu_\mu f\,d\lambda
       =-\int_F\widetilde{Gf}\,d\mu
                 \qquad(f\in L^\infty(F,\lambda)).
\end{equation}
Its canonical representative is finite and energy-quasicontinuous
outside a polar set, $T_ku_\mu\in H$, and $D(u_\mu S)=\mu$ on every
original compact scalar test. The map is linear, order reversing on
positive data, and
\begin{align}
 \|u_\mu-u_\nu\|_{L^1(\lambda)}
     &\le2\|\mu-\nu\|_{\rm TV},\label{eq:linear-l1-stability}\\
 \mathcal E(T_k(u_\mu-u_\nu))
     &\le k\|\mu-\nu\|_{\rm TV},\label{eq:linear-clipping-bound}\\
 c_V(\{|\widetilde u_\mu-\widetilde u_\nu|>k\})
     &\le k^{-1}\|\mu-\nu\|_{\rm TV}.
                         \label{eq:linear-capacity-bound}
\end{align}
\end{theorem}

\begin{proof}
\textbf{Step 1: finite-energy data.}
A finite diffuse measure is a smooth measure for the regular transient
form. The finite-energy nest theorem
\cite[Lemma 2.1]{KRStructure} gives increasing Borel sets $F_j$ with
$|\mu|(F\setminus\bigcup F_j)=0$ such that
$\mathbf1_{F_j}|\mu|$ has finite energy integral. Put
$\mu_j=\mathbf1_{F_j}\mu$ and $u_j=-G\mu_j$.

For a signed datum $\sigma$ with finite-energy variation, order gives
$|G\sigma|\le G|\sigma|$. Pair with the torsion vector $G1\le2$ to
obtain $\|G\sigma\|_1\le2\|\sigma\|_{\rm TV}$. The normal
contraction $s-2T_k(s)$ gives
$\mathcal E(T_kG\sigma)\le\mathcal E(G\sigma,T_kG\sigma)
\le k\|\sigma\|_{\rm TV}$. Testing capacity with
$|T_kG\sigma|/k$ gives the analogous capacity estimate.
Apply these estimates to $\mu_j-\mu_l$.

\textbf{Step 2: take the canonical limit.}
The $u_j$ are Cauchy in $L^1(\lambda)$ and in energy capacity. A
summable-capacity subsequence, together with the quasicontinuity of
its terms, has a finite energy-quasicontinuous limit off a polar set.
For fixed $k$, the $T_ku_j$ are bounded in $H$, and
$\|T_ku_j-T_ku\|_2^2\le2k\|u_j-u\|_1$. Weak compactness and
coercivity give $T_ku\in H$ and the clipping bound. The canonical
representatives agree by capacity convergence. Finite-energy symmetry
passes to \MathEqref{eq:transposition}, using total-variation
convergence of $\mu_j$ and boundedness of $Gf$.
Bounded measurable functions separate $L^1$ classes, so the
transposition identity gives uniqueness and linearity.

\textbf{Step 3: identify every original test.}
Coefficient domination makes $uS$ defined on all compact original
coefficients. For a compact original smooth $\phi$,
\MathEqref{eq:original-form-coefficients} gives
$G(b_\phi)=-\phi$. Use $f=b_\phi$ in
\MathEqref{eq:transposition} to obtain
\[
 \langle D(uS),\phi\rangle
     =\int_Fu b_\phi\,d\lambda=\int_F\phi\,d\mu.
\]
A common finite-energy exhaustion of $|\mu|+|\nu|$ passes the three
finite-energy difference estimates to the limit. This proves the
bounds and the order assertion.
\end{proof}

\subsection{Closed-form identities used for signed data}

For $v\in H$, its original core gradients close to a Hilbert gradient
module, and its energy measure $\Gamma_S(v)$ is the total-variation
limit of the core energy measures. In particular,
\begin{equation}\label{eq:energy-variation-norm}
 \|\Gamma_S(v)-\Gamma_S(w)\|_{\rm TV}
 \le\mathcal E(v-w)^{1/2}
       \bigl(\mathcal E(v)^{1/2}+\mathcal E(w)^{1/2}\bigr).
\end{equation}
The closed chain and product rules give
\begin{equation}\label{eq:hilbert-clipping}
 dT_kv=\mathbf1_{\{|\widetilde v|<k\}}\,dv,
 \qquad \Gamma_S(v)\{\widetilde v=t\}=0.
\end{equation}
For $h_t(s)=(s-t)^+-(-t)^+$, the finite-energy level measure is
\begin{equation}\label{eq:hilbert-level}
 L_t(v;S)=D(h_t(v)S)-\mathbf1_{\{\widetilde v>t\}}D(vS).
\end{equation}
When $|D(vS)|$ has finite energy integral, this is positive, diffuse
and carried by $\{\widetilde v=t\}$. These original-current facts
are proved in \cite[Sections 41--42]{Energy}. The scalar chain is
obtained by closing the original identity against tests
$\phi\Phi'(v)$; one-sided smooth hinges prove
\MathEqref{eq:hilbert-clipping} and the level assertion.

For a positive diffuse datum $\tau$, the potential $w=G\tau$ has
strongly convergent bounded clippings of its increasing finite-energy
nest potentials. Their energy identities give
\begin{equation}\label{eq:positive-band}
 \mathcal E(T_{2M}w-T_Mw)
                  \le M\tau\{w>M\}.
\end{equation}
The compatible energy measures of the clippings satisfy
\begin{equation}\label{eq:compatible-clippings}
 \Gamma_S(T_ku)
      =\mathbf1_{\{|\widetilde u|<k\}}\Gamma_S(T_lu),\qquad l>k.
\end{equation}
For any diffuse datum the Kato inequality, obtained from smooth
one-sided hinges in the transposition approximation, defines a positive
Radon measure
$L_t(u;S)=D(h_t(u)S)-\mathbf1_{\{\widetilde u>t\}}\mu$.
Its exact level carrier for signed data is proved below.

We use one elementary continuity fact. If $\nu$ is finite and
energy diffuse, then $c_V(E_j)\to0$ implies $|\nu|(E_j)\to0$.
Otherwise a subsequence with summable capacities and variations bounded
below would have a polar limsup with positive variation. Consequently,
capacity convergence passes bounded continuous scalar compositions in
$L^1(|\nu|)$.

\subsection{Signed clippings, chain rules and level measures}

Fix a finite signed energy-diffuse datum $\mu$ and write
$B_\mu=\|\mu\|_{\rm TV}$. Every occurrence of $u$ below means its
canonical representative when used as a measure multiplier.

\begin{lemma}[Bounded energy tests and signed Radon functionals]
\label[lemma]{lem:r11-bounded-H-tests-diffuse}
Every bounded $z\in H$ has original compact smooth core
approximations converging strongly in $H$, uniformly bounded
in absolute value, and along a subsequence quasi-everywhere.
If a finite signed Radon measure $\nu$ on $F$ has a continuous
core functional for the energy norm, then $|\nu|$ charges no
energy-polar set.  If $h\in H$ and $D(hS)=\nu$ on every original
compact smooth test, then
\begin{equation}\label{eq:signed-sd1}
       \mathcal E(h,z)=-\int\widetilde z\,d\nu
                     \qquad(z\in H\cap L^\infty),
       \qquad h=u_\nu .
\end{equation}


\end{lemma}

\begin{proof}
For bounded $z$, start with any original core sequence $z_j\to z$
in $H$ and compose it with one fixed smooth normal clamp that
fixes a neighborhood of the range of $z$ and is bounded globally.
The gradient-module chain rule proves strong convergence of these
compositions to $z$.  They are still original compact smooth
functions because the clamp fixes zero.  A summable-capacity
subsequence converges quasi-everywhere.

Let $K$ be an energy-polar compact set and choose an ordinary
relatively open neighborhood $G$ with compact closure such that
$|\nu|(G\setminus K)<\epsilon$.
The regular compact-core capacity formula, coercivity, a smooth
scalar clamp, and a fixed original cutoff supported in $G$
give core functions $0\le a_j\le1$, equal to one near $K$,
supported in $G$, with $\mathcal E(a_j)\to0$.
Multiplication by that cutoff preserves the energy convergence
by \MathEqref{eq:cutoff-energy}.  Hence $\int a_j\,d\nu\to0$, while
\[
       \left|\nu(K)-\int a_j\,d\nu\right|
                            \le |\nu|(G\setminus K)<\epsilon.
\]
Let $\epsilon\downarrow0$.  The same argument applies to every
polar compact subset of the positive and negative Hahn sets.
Radon inner regularity of the two variations therefore proves
that $|\nu|$ vanishes on every polar Borel set.

For bounded $z\in H$ use the uniformly bounded core
approximation just constructed.  Energy continuity passes its
left side in the original identity; diffuseness and dominated
convergence pass the right side.  This proves \MathEqref{eq:signed-sd1}'s first
assertion.  For $f\in L^\infty(\lambda)$ take $z=Gf$, which
is bounded by the torsion estimate.  Symmetry gives
\[
 \int h f\,d\lambda=\mathcal E(h,Gf)
                         =-\int\widetilde{Gf}\,d\nu.
\]
Since $h\in L^2(\lambda)\subset L^1(\lambda)$, transposition
uniqueness gives $h=u_\nu$.
\end{proof}

The usual product rule extends to bounded members of $H$:
$d(z_1z_2)=z_1dz_2+z_2dz_1$.  Indeed the preceding uniformly bounded
core approximations and the original gradient module pass this
identity by dominated convergence against the fixed limiting
energy measures.  This fact will be used for a product of two
bounded energy functions, not only a smooth cutoff.

\begin{theorem}[Strong clipping convergence for signed nest data]
\label[theorem]{thm:r11-signed-strong-truncations}
Let $\mu_j=\mathbf1_{F_j}\mu$ be the finite-energy nest
restrictions from the finite-energy nest in the proof of \cref{thm:linear-transposition}, and put $u_j=-G\mu_j$.
For every $k>0$,
\begin{equation}\label{eq:signed-sd2}
                          T_ku_j\longrightarrow T_ku_\mu
                                      \quad\hbox{strongly in }H.
\end{equation}

\end{theorem}

\begin{proof}
Let $w=-u_{|\mu|}=G|\mu|\ge0$ be the positive-data Green
potential, with its canonical representative.  Positivity and
order of the transposition map give
$|u_j|,|u_\mu|\le w$ quasi-everywhere for all $j$.
Put $a_M=|\mu|\{w>M\}$, so $a_M\to0$, and define
\[
          \zeta_M=\frac{T_{2M}w-T_Mw}{M}\in H,\qquad
          0\le\zeta_M\le1.
\]
It equals one on $\{w\ge2M\}$.  The positive-data band estimate \MathEqref{eq:positive-band} and the closed clipping rule give
\begin{equation}\label{eq:signed-sd3}
 \mathcal E(\zeta_M)\le\frac{a_M}{M},\qquad
 d\zeta_M=\frac1M\mathbf1_{\{M<w<2M\}}\,dw.
\end{equation}
Here $dw$ on the printed bounded band means the compatible
clipped gradient $d(T_{2M+1}w)$ there; no globally Radon
first-order current of the unbounded $w$ is required.
Test the finite-energy equation for $u_j$ with
$\zeta_M^2T_ku_j\in H\cap L^\infty$.
The original closed gradient product and clipping rules give
\[
 \int\zeta_M^2\,d\Gamma_S(T_ku_j)
  =-\int\zeta_M^2 T_k\widetilde u_j\,d\mu_j
                  -2\int\zeta_M T_k\widetilde u_j\,
                                      d\Gamma_S(u_j,\zeta_M).
\]
The first term in absolute value is at most $k a_M$.
The cross energy is carried where $M<w<2M$, and there
$|u_j|<2M$, so its $u_j$ factor may be clipped at $2M$.
Cauchy--Schwarz and \MathEqref{eq:linear-clipping-bound} therefore prove
\begin{equation}\label{eq:signed-sd4}
 \sup_j\Gamma_S(T_ku_j)\{w\ge2M\}
   \le k a_M+2k\sqrt{\mathcal E(\zeta_M)
                                 \sup_j\mathcal E(T_{2M}u_j)}
   \le k a_M+2k\sqrt{2B_\mu a_M}\longrightarrow0.
\end{equation}
This is the required uniform control of energy in the region
where positive and negative potentials could be large and cancel.

On $A_M=\{w\le M\}$ one has $|u_j-u_l|\le2M$.
The first-order clipping rule and \MathEqref{eq:linear-clipping-bound}, applied to the
finite-energy datum $\mu_j-\mu_l$, give
\begin{equation}\label{eq:signed-sd5}
 \|\mathbf1_{A_M}(du_j-du_l)\|^2
      \le\mathcal E(T_{2M+1}(u_j-u_l))
      \le(2M+1)\|\mu_j-\mu_l\|_{\rm TV}\longrightarrow0.
\end{equation}
Thus the localized gradients have a fixed limit in the
original Hilbert gradient module.  Its energy measure is
absolutely continuous with respect to $\lambda$.
For a smooth scalar $P$ with $P(0)=0$ and compactly supported
first derivative, bounded convergence against this fixed
limiting module measure proves that
$\mathbf1_{A_M}d(P(u_j))$ converges strongly.
Here $u_j\to u$ in $L^1(\lambda)$, and hence in measure;
continuity and boundedness of $P'$ supply the required
convergence of its scalar multipliers.
The energy of $P(u_j)$ outside $A_M$ is bounded by a fixed
multiple of the tail of a clipping $T_Lu_j$, where $L$
contains the support of $P'$.  Formula \MathEqref{eq:signed-sd4} makes it small
uniformly in $j$ as $M\to\infty$.
It follows that $P(u_j)$ is Cauchy in $H$; its $L^2$ limit
is $P(u)$, by boundedness and the $L^1$ convergence of $u_j$.

Finally approximate $T_k$ by smooth normal clamps
$P_\epsilon$ with uniform scalar error at most $\epsilon$,
whose derivatives differ from that of $T_k$ only where
$k-\epsilon<|s|<k+\epsilon$.  For $0<\epsilon<k/2$, the
finite-energy band formula and the source equation give
\begin{equation}\label{eq:signed-sd6}
 \begin{split}
 \mathcal E(T_ku_j-P_\epsilon(u_j))
  &\le\mathcal E(T_{k+\epsilon}u_j-T_{k-\epsilon}u_j)\\
  &=-\int(T_{k+\epsilon}\widetilde u_j
                 -T_{k-\epsilon}\widetilde u_j)\,d\mu_j
       \le2\epsilon B_\mu.
 \end{split}
\end{equation}
The smooth compositions are strongly convergent by the
preceding paragraph.  First let $j,l\to\infty$ and then
$\epsilon\downarrow0$ in the triangle inequality.
This makes $T_ku_j$ Cauchy in $H$.  Its already known
$L^2$ limit is $T_ku$, proving \MathEqref{eq:signed-sd2}.
\end{proof}

\begin{theorem}[Full signed renormalized chain and energy distribution]
\label[theorem]{thm:r11-signed-renormalized-chain}
For $u=u_\mu$, the increasing compatible energy measure
$\Gamma_S(u)=\lim_k\Gamma_S(T_ku)$ is sigma-finite, finite
on every bounded value strip, and zero on every exact scalar
level.  If $\Phi\in C^2(\mathbf R)$ has bounded first derivative
and compactly supported second derivative, then
\begin{equation}\label{eq:signed-sd7}
 D(\Phi(u)S)=\Phi'(\widetilde u)\mu
                           +\Phi''(\widetilde u)\Gamma_S(u)
\end{equation}
on all original compact smooth tests.  Every term is coefficient-
defined, and the right side is a finite signed Radon measure.
If in addition $\Phi'$ has compact support, then
$\Phi(u)-\Phi(0)\in H$.
For every $0\le a<b$,
\begin{equation}\label{eq:signed-sd8}
 \mathcal E(T_bu-T_au)
                   =-\int(T_b\widetilde u-T_a\widetilde u)\,d\mu.
\end{equation}
The scalar energy image has density
\begin{equation}\label{eq:signed-sd9}
 \widetilde u_*\Gamma_S(u)=k_\mu(t)\,dt,\qquad
 k_\mu(t)=
 \begin{cases}
       \mu\{\widetilde u<t\},&t<0,\\
       -\mu\{\widetilde u>t\},&t>0.
 \end{cases}
\end{equation}
This density is nonnegative and bounded by $B_\mu$ almost everywhere;
its value at zero is immaterial.  In particular all bounded-value
weighted energy measures used here are finite Radon measures.
\end{theorem}

\begin{proof}
Compatibility and level zero follow from the $H$ clipping
identities exactly as in \MathEqref{eq:compatible-clippings} and \MathEqref{eq:hilbert-clipping}, without any sign
hypothesis on $u$.  Choose $M$ containing the support of
$\Phi''$ in its interior.  The finite-energy chain for $u_j$
reads
\[
 D(\Phi(u_j)S)=\Phi'(\widetilde u_j)\mu_j
                         +\Phi''(\widetilde u_j)\Gamma_S(T_Mu_j).
\]
The left side passes because $\Phi$ is Lipschitz and
$u_j\to u$ in $L^1(\lambda)$.  The source term passes in
variation by the capacity convergence of $u_j$
to $u$, diffuseness of $|\mu|$, bounded continuity of $\Phi'$, and
total-variation convergence of $\mu_j$ to $\mu$.
Strong convergence in \MathEqref{eq:signed-sd2}, the energy variation estimate
\MathEqref{eq:energy-variation-norm}, and bounded convergence against the fixed limiting
energy measure pass the last term.  This proves \MathEqref{eq:signed-sd7}.
For merely $C^2$ scalar functions, ordinary smoothing gives
the same limit, uniformly in the first two derivatives on
the bounded active range and uniformly in scalar values
up to the original Lipschitz error.
If $\Phi'$ is compactly supported, its normalized composition
is a smooth function of one bounded clipping and lies in $H$.

The same strong clipping convergence passes the finite-energy
band identity, proving \MathEqref{eq:signed-sd8}.  Apply the argument separately
to the bounded one-sided bands
\[
 B^+_{a,b}(s)=\min((s-a)^+,b-a),\qquad
 B^-_{a,b}(s)=-B^+_{a,b}(-s).
\]
Their derivatives are the indicators of $(a,b)$ and
$(-b,-a)$, respectively.  Composing a strongly convergent
clipping with these finite hinge combinations is strongly
continuous: the closed gradient formula applies, and the
limiting $H$ energy is zero on their finitely many corner
levels by \MathEqref{eq:hilbert-clipping}.  Hence
\[
 \begin{split}
 \Gamma_S(u)\{a<u<b\}&=-\int_a^b\mu\{u>t\}\,dt,\\
 \Gamma_S(u)\{-b<u<-a\}&=\int_a^b\mu\{u<-t\}\,dt.
 \end{split}
\]
There is no energy mass at zero or at the interval endpoints.
These formulas prove \MathEqref{eq:signed-sd9}.  Nonnegativity follows because
it is the density of a positive measure; the bound follows
from total variation of $\mu$.
\end{proof}

\begin{proposition}[Finite global level measures and clipping curvatures]
\label[proposition]{prop:r11-signed-global-curvature-truncation}
The Kato measures $L_t(u;S)$ are finite for every real $t$,
with
\begin{equation}\label{eq:signed-sd10}
 \begin{aligned}
 L_t(u;S)(F)&\le |\mu|\{u\le t\}&& (t<0),\\
 L_t(u;S)(F)&\le |\mu|\{u>t\}&& (t>0),\\
 L_0(u;S)(F)&\le B_\mu.
 \end{aligned}
\end{equation}
For $M>0$, the actual signed Radon curvature
$\nu_M=D(T_Mu\,S)$ satisfies
\begin{equation}\label{eq:signed-sd11}
 \begin{split}
 \nu_M&=\mathbf1_{\{-M<u\le M\}}\mu+L_{-M}(u;S)-L_M(u;S),\\
 \|\nu_M\|_{\rm TV}&\le B_\mu,\qquad
 \|\nu_M-\mu\|_{\rm TV}
                    \le2|\mu|\{|u|\ge M\}\longrightarrow0.
 \end{split}
\end{equation}
Both $\nu_M$ and every $L_t(u;S)$ are energy diffuse.

\end{proposition}

\begin{proof}
Apply \MathEqref{eq:signed-sd7} to the smooth one-sided hinges.  Their positive
defects are
$\Phi_\epsilon''(u-t)\Gamma_S(u)$, supported in the scalar
strip $t<u<t+\epsilon$.  Formula \MathEqref{eq:signed-sd9}, the nonnegative
unit-integral second derivative, and its bounded density
show that their masses are at most $B_\mu$.
For $t>0$ the sharper bound is $|\mu|\{u>t\}$; for
$t<0$ and sufficiently small $\epsilon$ it is
$|\mu|\{u<t+\epsilon\}$.
They converge on original compact tests to $L_t(u;S)$.
Positivity, compact smooth exhaustion and then
$\epsilon\downarrow0$ prove \MathEqref{eq:signed-sd10}.

The scalar identity $T_Ms=h_{-M}(s)-h_M(s)$ gives the
first formula of \MathEqref{eq:signed-sd11}.  The three source regions
$\{u\le-M\}$, $\{-M<u\le M\}$ and $\{u>M\}$
partition the finite-value locus.  Triangle inequalities
and \MathEqref{eq:signed-sd10} therefore give both variation estimates in
\MathEqref{eq:signed-sd11}.  Since $T_Mu\in H$, its finite Radon curvature
is diffuse by \cref{lem:r11-bounded-H-tests-diffuse}.

For any fixed $t$ and $M>|t|$, the exact scalar identity
\begin{equation}\label{eq:signed-sd12}
                    h_t(T_Mu)=h_t(u)-h_M(u)
\end{equation}
holds on the finite-value locus and hence in the coefficient-defined
$L^1$ sense.  Its left side belongs to $H$.  Its curvature is finite Radon
by \MathEqref{eq:signed-sd10} and the Kato decompositions, hence diffuse by
the same lemma.  Thus
\[
 L_t(u;S)-L_M(u;S)
       =D(h_t(T_Mu)S)-\mathbf1_{\{t<u\le M\}}\mu
\]
is a diffuse finite signed measure.  On any polar Borel
set, the mass of $L_t$ is therefore at most $L_M(F)$,
which tends to zero by \MathEqref{eq:signed-sd10}.  This proves diffuseness
of $L_t$ without assuming it passes through an arbitrary
weak measure limit.
\end{proof}

\begin{lemma}[Finite signed Radon data in the $H$ level theorem]
\label[lemma]{lem:r11-H-level-carrier-finite-signed}
Let $v\in H$ and suppose $D(vS)=\sigma$ is a finite signed
Radon measure on the original tests.  Then the conclusion
\MathEqref{eq:hilbert-level} and its exact carrier assertion holds, with a finite diffuse positive level
measure, without assuming finite energy integral of
$|\sigma|$.
\end{lemma}

\begin{proof}
The datum is diffuse and the current equation extends to
all bounded $H$ tests by Lemma
\ref{lem:r11-bounded-H-tests-diffuse}.
Consequently the finite-energy scalar chain proof applies:
its test $\psi\Phi'(v)$ is bounded and lies in $H$.
The one-sided band tests are also bounded $H$ vectors and
give \MathEqref{eq:signed-sd9} with $v,\sigma$ in place of $u,\mu$.
Thus the smoothed hinge defects have uniformly bounded
total mass, and their positive limit $L_t(v;S)$ is finite.
The signed measure
$D(h_t(v)S)=\mathbf1_{\{v>t\}}\sigma+L_t(v;S)$ is finite
and comes from an $H$ vector.  It is therefore diffuse by
the preceding lemma, as is $L_t$.

The original current identity for this hinge source now
extends to every bounded $H$ test.  Take
$z=\psi g(v)$ with $\psi$ an original compact nonnegative
smooth function and $g$ smooth bounded, with bounded
derivative, nonnegative and zero near $t$.
Strong convergence of the smooth hinges in $H$ and bounded
convergence against the finite measure $|\sigma|$ pass
\[
 \int z\,dL_{t,\epsilon}
     =-\mathcal E(h_{t,\epsilon}(v),z)
              -\int z\Phi_\epsilon'(v-t)\,d\sigma
\]
to $\int z\,dL_t$.
Its left side is zero for small $\epsilon$ because
$g\Phi_\epsilon''=0$.  The same countable family of
nonnegative tests as in the finite-energy level-carrier proof cited above shows that
$L_t$ is carried by $\{\widetilde v=t\}$.
\end{proof}

\begin{theorem}[Exact level carriers and coarea for all signed diffuse data]
\label[theorem]{thm:r11-signed-diffuse-exact-level}
For every finite signed energy-diffuse datum and every real $t$,
\begin{equation}\label{eq:signed-sd13}
                    \mathbf1_{\{\widetilde u\ne t\}}L_t(u;S)=0.
\end{equation}
For every bounded compactly supported Borel scalar $\theta$,
\begin{equation}\label{eq:signed-sd14}
                \int\theta(t)L_t(u;S)\,dt
                         =\theta(\widetilde u)\Gamma_S(u)
\end{equation}
as finite signed Radon measures.  The hard-curvature
variation bounds \MathEqref{eq:signed-sd11} hold on the complete original
space and every actual moment refinement.
\end{theorem}

\begin{proof}
First identify $\nu_M=\mu$ on $A_M=\{|u|<M\}$.
Choose a smooth $\Phi$ with $\Phi'$ zero outside $[-M,M]$
and strictly positive on $(-M,M)$.
The identities $\Phi(T_Mu)=\Phi(u)$ and \MathEqref{eq:compatible-clippings} make their
weighted energies agree.  Compare \MathEqref{eq:signed-sd7} with the bounded
$H$ chain for $T_Mu$ and its finite signed datum $\nu_M$.
The latter chain is valid by the bounded-test extension,
so $\Phi'(u)(\nu_M-\mu)=0$, proving the assertion on $A_M$.

Fix $t$, choose $M>|t|$, and apply the preceding $H$ level
lemma to $v=T_Mu$.  The level measure $L_t(v;S)$ is carried
by $\{v=t\}=\{u=t\}$.
Compare the curvature of the two sides of \MathEqref{eq:signed-sd12}, and then
restrict their actual Radon measures to $A_M$.
The just proved identity $\nu_M=\mu$ on $A_M$ cancels the
source terms and gives
\[
     \mathbf1_{A_M}(L_t(u;S)-L_M(u;S))
                         =\mathbf1_{A_M}L_t(T_Mu;S).
\]
It follows that
$L_t(u;S)(A_M\cap\{u\ne t\})\le L_M(u;S)(F)\to0$.
The sets $A_M$ exhaust the finite-value locus, and $L_t$ is
diffuse, so it gives no mass to the omitted polar infinite-
value set.  This proves \MathEqref{eq:signed-sd13} by monotone measure convergence.

For smooth compactly supported $\theta$, take
$\Phi(s)=\int\theta(t)(s-t)^+dt$.
It has bounded first derivative and compactly supported
second derivative $\theta$, so \MathEqref{eq:signed-sd7} and Fubini prove
\MathEqref{eq:signed-sd14} on every original test.
The positive kernels $L_t$, their global bound by $B_\mu$,
and their Borel original-test pairings extend the equality
to nonnegative Borel $\theta$; subtract positive and negative
parts. The right side is finite by \MathEqref{eq:signed-sd9}.\qedhere
\end{proof}

\begin{definition}[The signed renormalized zero-boundary class]
\label[definition]{def:r11-signed-renormalized-boundary}
For the specified finite signed diffuse datum, a candidate is
$v\in L^1(F,\lambda)$ with a finite energy-quasicontinuous
representative and $T_kv\in H$ for every $k>0$.
Use the compatible sigma-finite energy measure of its
clippings.  Require the original current identity
\begin{equation}\label{eq:signed-sd15}
 D(\Phi(v)S)=\Phi'(\widetilde v)\mu
                     +\Phi''(\widetilde v)\Gamma_S(v)
\end{equation}
for every $\Phi\in C^2(\mathbf R)$ with compactly supported
first derivative, and the radiation condition
$k^{-1}\mathcal E(T_{2k}v-T_kv)\to0$.
These conditions specify the boundary and solution class
without using the transposition identity in their definition.
\end{definition}

\begin{theorem}[The complete finite-diffuse relative renormalized Dirichlet problem]
\label[theorem]{thm:r11-signed-renormalized-Dirichlet}
Every finite signed energy-diffuse Radon datum has exactly one
candidate in \cref{def:r11-signed-renormalized-boundary}.
It is the transposition solution $u_\mu$, and all its original
current, exact-level, coarea and global curvature-variation
conclusions above hold.  The construction and solution class
are independent of the chosen faithful reference measure.
\end{theorem}

\begin{proof}
Existence follows from \MathEqref{eq:signed-sd7}, the truncation membership,
quasicontinuity, and the band identity \MathEqref{eq:signed-sd8}, which gives $k^{-1}\mathcal E(T_{2k}u-T_ku)\le|\mu|\{|u|>k\}\to0$.  For uniqueness choose smooth odd scalar
clamps $\Phi_k$ fixing the interval $[-k,k]$, constant
outside $[-2k,2k]$, with $\Phi_k(0)=0$,
$0\le\Phi_k'\le1$ and
$|\Phi_k''|\le C/k$ on the two transition bands.
Then $|\Phi_k(s)|\le|s|$ and $\Phi_k(s)\to s$.
For a candidate $v$, its prescribed chain gives
\begin{equation}\label{eq:signed-sd16}
 \begin{split}
 \nu_k:=D(\Phi_k(v)S)
       &=\Phi_k'(v)\mu+\Phi_k''(v)\Gamma_S(v),\\
 \|\nu_k-\mu\|_{\rm TV}
       &\le |\mu|\{|v|>k\}
                +\frac Ck\mathcal E(T_{2k}v-T_kv)\to0.
 \end{split}
\end{equation}
The measures $\nu_k$ are finite signed Radon measures.
The potential $\Phi_k(v)=\Phi_k(T_{2k}v)$ belongs to $H$,
so \cref{lem:r11-bounded-H-tests-diffuse} proves
both diffuseness of $\nu_k$ and
$\Phi_k(v)=u_{\nu_k}$.  
The established total-variation stability gives
$\|\Phi_k(v)-u_\mu\|_1\le2\|\nu_k-\mu\|_{\rm TV}\to0$.
Dominated convergence also gives $\Phi_k(v)\to v$ in
$L^1(\lambda)$.  Thus $v=u_\mu$ almost everywhere.
Equality of all clipped $H$ vectors identifies the canonical
representatives quasi-everywhere.  Reference independence follows because the original core energy
$\mathcal E(f,g)=\int d'f\wedge d''g\wedge S$ is independent of the
faithful measure. Coercivity identifies each realization with the same
Hilbert completion. Its compact-core zero-capacity, polar sets and
canonical representatives coincide. A common finite-energy nest then
gives the same capacity limit of the same Riesz potentials.
\end{proof}


\subsection{Stability under signed measure convergence}

\begin{theorem}[Strong clipping and smooth-curvature stability]
\label[theorem]{thm:r11-signed-data-strong-clipping}
Suppose $\|\mu_j-\mu\|_{\rm TV}\to0$, and write
$u_j=u_{\mu_j}$, $u=u_\mu$ for the unique intrinsic
renormalized solutions of $D(u_jS)=\mu_j$, $D(uS)=\mu$.
Then, for every fixed $k>0$,
\begin{equation}\label{eq:stability-ds1}
                    \|T_ku_j-T_ku\|_H\longrightarrow0.
\end{equation}
In particular, for any fixed $k,l>0$, the finite original
energy measures converge in global total variation:
\begin{equation}\label{eq:stability-ds2}
 \Gamma_S(T_ku_j,T_lu_j)
               \longrightarrow\Gamma_S(T_ku,T_lu).
\end{equation}
If $\Phi\in C^2(\mathbf R)$ has bounded first derivative
and compactly supported second derivative, then
\begin{equation}\label{eq:stability-ds3}
           D(\Phi(u_j)S)\longrightarrow D(\Phi(u)S)
                       \quad\hbox{in global total variation}.
\end{equation}
If, more restrictively, $\Phi'$ has compact support, then
$\Phi(u_j)-\Phi(0)\to\Phi(u)-\Phi(0)$ strongly in $H$.
For this latter class both original first-order currents
$d'(\Phi(u_j)S)$ and $d''(\Phi(u_j)S)$ converge in local
coefficient total variation, against every original compact
complementary smooth test and every actual moment refinement.
The zero-order multipliers $\Phi(u_j)S$ converge in that same
local coefficient sense under the hypotheses of \MathEqref{eq:stability-ds3}.
\end{theorem}

\begin{proof}
The finite-datum estimates give $u_j\to u$ in energy capacity
and in $L^1(\lambda)$.  The clipped functions converge in
$L^2(\lambda)$: they are bounded by $k$, and $\lambda(F)<\infty$
because the part domain is relatively compact in the neighborhood
where $\lambda$ is Radon.  Also
$\mathcal E(T_ku_j)\leq k\|\mu_j\|_{\rm TV}$, so these
vectors are bounded in $H$.  Every weakly convergent subsequence
therefore has the unique $L^2$ representative $T_ku$.
Consequently the whole sequence converges weakly in $H$.

Take $a=0,b=k$ in \MathEqref{eq:signed-sd8}.  The exact norm identity is
\begin{equation}\label{eq:stability-ds4}
 \mathcal E(T_ku_j)=-\int_F T_k\widetilde u_j\,d\mu_j,
 \qquad
 \mathcal E(T_ku)=-\int_F T_k\widetilde u\,d\mu.
\end{equation}
The difference between the first datum integral and the second
is bounded in absolute value by
\begin{equation}\label{eq:stability-ds5}
 k\|\mu_j-\mu\|_{\rm TV}
       +\int_F|T_k\widetilde u_j-T_k\widetilde u|\,d|\mu|.
\end{equation}
The second term tends to zero by energy-capacity convergence,
finite diffuseness of $|\mu|$, and boundedness of the clippings.
Thus the norms in \MathEqref{eq:stability-ds4} converge.  Weak convergence and convergence
of norms in the Hilbert space prove \MathEqref{eq:stability-ds1}.  The exact band identity is what turns weak clipping convergence
into norm convergence.

The original closed gradient module gives, for $f_j,g_j,f,g\in H$,
\begin{equation}\label{eq:stability-ds6}
 \|\Gamma_S(f_j,g_j)-\Gamma_S(f,g)\|_{\rm TV}
  \leq \|f_j-f\|_H\|g_j\|_H
                         +\|f\|_H\|g_j-g\|_H.
\end{equation}
This is measure Cauchy--Schwarz after adding and subtracting
$\Gamma_S(f,g_j)$.  Applying \MathEqref{eq:stability-ds1} proves \MathEqref{eq:stability-ds2}.

For \MathEqref{eq:stability-ds3}, choose $M$ larger than the support radius of $\Phi''$.
The all-original chain identity reads
\begin{equation}\label{eq:stability-ds7}
 D(\Phi(u_j)S)=\Phi'(\widetilde u_j)\mu_j
       +\Phi''(\widetilde u_j)\Gamma_S(T_Mu_j).
\end{equation}
The first term converges in variation: its changing-datum error
is at most $\|\Phi'\|_\infty\|\mu_j-\mu\|_{\rm TV}$,
and its fixed-datum multiplier error tends to zero by finite
diffuseness and capacity convergence.  For the second term use
\MathEqref{eq:stability-ds2} and then bounded convergence in measure against the fixed
finite diffuse measure $\Gamma_S(T_Mu)$.  This gives its variation
convergence as well and proves \MathEqref{eq:stability-ds3}.  The zero-order multiplier
claim follows from the Lipschitz bound for $\Phi$, $L^1(\lambda)$
stability and the faithful original coefficient domination.

If $\Phi'$ has compact support, its normalized composition is
a smooth function of one sufficiently large clipping.  The
closed gradient formula and \MathEqref{eq:stability-ds1} give strong convergence:
the changing-gradient term is bounded by
$\|\Phi'\|_\infty\|T_Mu_j-T_Mu\|_H$, and the
changing scalar factor tends to zero in the fixed finite
energy measure.  Its $L^2$ representative converges by the
same bounded composition argument.  Finally, on a fixed
original compact stage, write any complementary one-form
test as a finite sum of bounded smooth coefficients times
differentials of original local moment functions.  Multiply
those moment functions by original compact cutoffs equal to
one on the test support.  The coefficient pairing of a
first-order derivative with each such term is an original
mixed gradient measure; \MathEqref{eq:stability-ds6}, or its one-variable
Cauchy--Schwarz bound, controls its variation by a fixed
constant times the $H$-norm of the changing potential.
This proves local coefficient-TV convergence of both first-order
currents.  The argument starts with an arbitrary original test
and is valid after any further actual moment refinement.

\end{proof}


% END sections/10_linear_dirichlet.tex

% BEGIN sections/11_bibliography.tex
\begin{thebibliography}{KR17}
\bibitem[APW26]{APW26}
O.~Amini, M.~Piquerez and L.~Wu,
\emph{On solutions to tropical Monge--Amp\`ere equations},
arXiv:2610.25602v1 (2026).
\url{https://arxiv.org/abs/2610.25602}.

\bibitem[BMPS26]{BMPS26}
A.~M.~Botero, E.~Mazzon and L.~Pille--Schneider,
\emph{Monge--Amp\`ere measures on balanced polyhedral spaces},
arXiv:2603.08664v1 (2026).
\url{https://arxiv.org/abs/2603.08664}.

\bibitem[Ceg98]{Cegrell98}
U.~Cegrell, \emph{Pluricomplex energy},
Acta Math. \textbf{180} (1998), 187--217.
\url{https://doi.org/10.1007/BF02392899}.

\bibitem[CLD12]{CLD12}
A.~Chambert--Loir and A.~Ducros,
\emph{Formes diff\'erentielles r\'eelles et courants sur les espaces de
Berkovich}, arXiv:1204.6277v1 (2012).
\url{https://arxiv.org/pdf/1204.6277v1}.

\bibitem[CLD25]{CLD25}
A.~Chambert--Loir and A.~Ducros,
\emph{Formes diff\'erentielles r\'eelles et courants sur les espaces de
Berkovich}, arXiv:1204.6277v2 (2025).
\url{https://arxiv.org/pdf/1204.6277v2}.

\bibitem[DE06]{DE06}
K.~Dabbek and F.~Elkhadhra,
\emph{Capacity associated to a positive closed current},
Documenta Math. \textbf{11} (2006), 469--486.
\url{https://doi.org/10.4171/DM/219}.

\bibitem[GR25]{GR25}
W.~Gubler and J.~Rabinoff,
\emph{Classically psh and pluriharmonic functions on Berkovich spaces},
arXiv:2506.13548v2 (2025).
\url{https://arxiv.org/pdf/2506.13548v2}.

\bibitem[KR15]{KR15}
T.~Klimsiak and A.~Rozkosz,
\emph{Renormalized solutions of semilinear equations involving measure
data and operator corresponding to Dirichlet form},
NoDEA Nonlinear Differential Equations Appl. \textbf{22} (2015),
1911--1934.
\url{https://doi.org/10.1007/s00030-015-0350-1}.

\bibitem[KR17]{KRStructure}
T.~Klimsiak and A.~Rozkosz,
\emph{On the structure of bounded smooth measures associated with a
quasi-regular Dirichlet form},
Bull. Polish Acad. Sci. Math. \textbf{65} (2017), 45--56;
arXiv:1410.4927v2, Lemma 2.1.
\url{https://doi.org/10.4064/ba8108-7-2017}.

\bibitem[Vu21]{Vu21}
D.-V.~Vu, \emph{Relative non-pluripolar product of currents},
Ann. Global Anal. Geom. \textbf{60} (2021), 269--311.
\url{https://doi.org/10.1007/s10455-021-09780-7}.

\bibitem[Wu26]{Wu26}
L.~Wu, \emph{On the local theory of non-archimedean psh functions},
arXiv:2609.27610v1 (2026).
\url{https://arxiv.org/abs/2609.27610}.

\bibitem[Cont]{Continuous}
\emph{Continuous CLD-positive potentials and mixed currents},
internal manuscript, version 2.6, 15 November 2026.
The supplied source is \nolinkurl{R_CLD_Regularization_v2.6.tex}.

\bibitem[Bound]{Bounded}
\emph{Locally bounded classical plurisubharmonic functions
and mixed currents}, internal manuscript, edited version 0.1,
10 November 2026.
The supplied source is \nolinkurl{B_Agent_Edited_v0.1.tex}.

\bibitem[Rel]{Relative}
\emph{Relative Bedford--Taylor theory on Berkovich spaces},
internal research manuscript, fourth record, 11 November 2026.
The bounded relative energy and original coefficient arguments are
retained in the accompanying complete source and proof archive.

\bibitem[Quot]{Quotient}
\emph{Bounded classical curvature kernels and arbitrary-background
potential quotients}, internal research manuscript, fifth record,
10 November 2026. See in particular Section 10.

\bibitem[Order]{Order}
\emph{The ordered potential quotient, delta tests, and relative
Bedford--Taylor theory}, internal research manuscript, sixth record,
10 November 2026. The supplied complete source is
\nolinkurl{NA_BT_R6_Delta_Capacity_Dirichlet_2026-11-10.tex};
its formal title is the one printed here.

\bibitem[Energy]{Energy}
\emph{Unbounded energy, capacity, and relative Dirichlet theory},
internal research manuscript, eleventh record, 11 November 2026.
The complete source and PDF are supplied as companions. Section and
parenthesized estimate references in this article refer to that
fixed version.

\bibitem[Dual]{Dual}
\emph{Closed bands, energy duals, and nonlinear source gaps},
internal research manuscript, twelfth record, 12 November 2026.
See Sections 12--13, 25 and 31 for the dual and probe constructions.

\bibitem[OrdComp]{Intrinsic}
\emph{Order-controlled intrinsic extensions and sufficient certificates
for non-Archimedean Bedford--Taylor actions}, internal proof note,
14 November 2026, Sections 2--5.
\end{thebibliography}

% END sections/11_bibliography.tex

\end{document}
